\documentclass[11pt]{article}
\usepackage[utf8]{inputenc}
\usepackage[T1]{fontenc}
\usepackage{amsmath,amssymb,amsthm,mathtools}
\usepackage{graphicx}
\usepackage{hyperref}
\usepackage{booktabs}
\usepackage{enumitem}
\usepackage[margin=1in]{geometry}
\graphicspath{{figures/}}
\setlength{\emergencystretch}{2em}

\newtheorem{theorem}{Theorem}[section]
\newtheorem{proposition}[theorem]{Proposition}
\newtheorem{lemma}[theorem]{Lemma}
\newtheorem{corollary}[theorem]{Corollary}
\theoremstyle{definition}
\newtheorem{definition}[theorem]{Definition}
\theoremstyle{remark}
\newtheorem{remark}[theorem]{Remark}
\newtheorem{example}[theorem]{Example}

\newcommand{\R}{\mathbb{R}}
\newcommand{\E}{\mathbb{E}}
\newcommand{\Prob}{\mathbb{P}}
\newcommand{\norm}[1]{\left\|#1\right\|}
\newcommand{\nop}[1]{\left\|#1\right\|_{\mathrm{op}}}
\newcommand{\nnuc}[1]{\left\|#1\right\|_{*}}
\newcommand{\nF}[1]{\left\|#1\right\|_{F}}
\newcommand{\inner}[2]{\left\langle #1, #2 \right\rangle}
\newcommand{\abs}[1]{\left|#1\right|}
\newcommand{\tr}{\operatorname{tr}}
\newcommand{\diag}{\operatorname{diag}}
\newcommand{\polar}{\operatorname{polar}}
\newcommand{\sharpop}{\operatorname{sharp}}
\newcommand{\srank}{\operatorname{srank}}
\newcommand{\nr}{\operatorname{nr}}
\newcommand{\sign}{\operatorname{sign}}
\newcommand{\Var}{\operatorname{Var}}
\newcommand{\MAD}{\operatorname{MAD}}
\newcommand{\Med}{\operatorname{Med}_\ast}
\newcommand{\Lop}{L_{\mathrm{op}}}
\newcommand{\muop}{\mu_{\mathrm{op}}}
\newcommand{\kop}{\kappa_{\mathrm{op}}}
\newcommand{\kF}{\kappa_{F}}
\newcommand{\cO}{\mathcal{O}}
\newcommand{\lmin}{\lambda_{\min}}
\newcommand{\lmax}{\lambda_{\max}}
\newcommand{\ones}{\mathbf{1}}

\title{Polar (Muon-type) steps on strongly convex matrix problems:\\ condition numbers, exact dynamics, and the effect of noise}
\author{Theory notes of the side study \texttt{muon/} --- project ``Accelerating Variance Reduction''\\ (Statistical Learning Theory, UC Berkeley, November 2026)}
\date{Version of 16 November 2026, revised after an adversarial review}

\begin{document}
\maketitle

\begin{abstract}
These notes push the theoretical side of the question ``what do the results of the main project say about Muon?''.
Muon moves a matrix parameter along the polar factor of (a momentum of) its gradient: steepest descent for the spectral
norm. The main project showed, for SAGA, that worst-case constants hide what happens along trajectories (notes B2 and
B3) and that no fixed metric improves the global condition number of logistic regression (note B1). A first version of
these notes derived the counterparts for the spectral geometry; it was then reviewed by six independent referee passes
(three on the mathematics, three on the literature, 16 November 2026). The mathematics survived with corrections;
the literature review showed that a large part of the dynamics on least squares and of the noise analysis was already
published in 2025--2026. This version separates what is known (attributed, with the statement numbers read by the
referees) from what is added. The additions are: (i) on conditioning, a frame bound refining the dimension barrier
$\kop\ge\min(K,d)$, with its rigidity, a sharp comparison with the Euclidean condition number through nuclear and stable
ranks, and a \emph{no-free-lunch principle valid for every norm}: for ridge-regularized generalized linear models whose
loss curvature vanishes along a sequence of parameters (logistic, softmax), no norm, spectral or other, has a better
worst-case condition number than the Frobenius norm; (ii) the reading of polar
steps as mirror \emph{subgradient} descent on the nuclear norm of the gradient (the nonsmooth case left out by the
theory of dual-space preconditioning), which makes the nuclear-scaled step a Polyak step and gives a geometric schedule
for every smooth strongly convex objective; (iii) on the least-squares model, whose dynamics are mostly published, a
rotation factor between the error and its gradient, an exact conservation law in dimension two, a counterexample to
the monotonicity of the spectral norm of the error in dimension three, and a sharp contraction bound after the first
step; (iv) on noise: stochastic polar steps with additive noise are stochastic mirror descent on a noise-smoothed
nuclear norm, hence, for lengths with $\sum_t\eta_t=\infty$ and $\sum_t\eta_t^2<\infty$, converge almost surely to a
point where zero is a \emph{nuclear median}, not the mean, of the stochastic gradient; for batch one on least squares
the limit is a weighted least-distances estimator; for skewed gradient noise with independent entries a formal
expansion gives an offset of order $1/b$ with explicit constants, smaller per entry than for sign descent by a
factor of order $1/d$; a finite-sum variance-reduced estimator has a fixed point without offset. Every statement
carries its status; ``not found'' means not found by the literature checks of 12 and 16 November 2026, and the
statements added during the revision, which were not searched separately, are listed in Appendix~\ref{app:review}.
\end{abstract}

\tableofcontents

\section{Setting, notation, status of the statements}\label{sec:setting}

\paragraph{Matrices and norms.} $W\in\R^{K\times d}$, $r=\min(K,d)$, $\inner{A}{B}=\tr(A^\top B)$. $\nF{\cdot}$ is the
Frobenius norm, $\nop{\cdot}$ the spectral norm, $\nnuc{\cdot}$ the nuclear norm (its dual:
$\abs{\inner AB}\le\nnuc A\nop B$). For $M=U\diag(s)V^\top$ (thin singular value decomposition restricted to the nonzero
singular values) we write
\[
\polar(M)=UV^\top,\qquad \sharpop(M)=\nnuc M\,\polar(M),\qquad \abs{M}=(M^\top M)^{1/2},
\]
$\srank(M)=\nF M^2/\nop M^2$ (stable rank) and $\nr(M)=\nnuc M^2/\nF M^2$ (nuclear rank, the term of Davis and
Drusvyatskiy); $\srank(M)\le\nr(M)$. The function $\psi(G)=\nnuc G$ is convex with subdifferential
$\partial\psi(G)=\{UV^\top+P_{U^\perp}CP_{V^\perp}:\nop C\le1\}$; $\polar(G)$ is its element of least Frobenius norm,
$\inner GO=\nnuc G$ for every $O\in\partial\psi(G)$, and $\nop O=1$ if $G\ne0$. With this convention $\polar(0)=0$:
statements about floors below need $G\ne0$, statements about upper bounds hold for every selection
$O\in\partial\psi(G)$. $\cO$ denotes the set of $K\times d$ matrices whose $r$ singular values are all equal to one.

\paragraph{Objective.} $F:\R^{K\times d}\to\R$ is twice continuously differentiable, with minimizer $W_\ast$,
$F_\ast=F(W_\ast)$, $G=\nabla F(W)$, $E=W-W_\ast$. The constants of $F$ in the spectral geometry are the best numbers
such that, for all $W$ and all directions $\Delta$,
\begin{equation}\label{eq:consts}
\muop\nop\Delta^2\ \le\ D^2F(W)[\Delta,\Delta]\ \le\ \Lop\nop\Delta^2,\qquad \kop=\Lop/\muop ,
\end{equation}
and $\mu_F$, $L_F$, $\kF=L_F/\mu_F$ denote the same constants for the Frobenius norm; for another norm $N$ we write
$\mu_N,L_N,\kappa_N$. Since $\nop\Delta\le\nF\Delta\le\sqrt r\nop\Delta$, one has $L_F\le\Lop\le rL_F$ and
$\mu_F\le\muop\le r\mu_F$.

\paragraph{Steps.} A \emph{polar step of length $\eta$} is $W^+=W-\eta O$ with $O\in\partial\psi(G)$ (in practice
$O=\polar(G)$); with a scheduled $\eta$ and a momentum of stochastic gradients in place of $G$ it is Muon. The
\emph{nuclear-scaled step} is $W^+=W-\frac{1}{L}\sharpop(G)$: the ``sharp operator'' of stochastic spectral descent
(Carlson, Cevher and Carin, AISTATS 2015), called PolarGrad by Lau, Long and Su (arXiv:2505.21799).

\paragraph{Two models.} (LS) the \emph{one-sided quadratic}
$F(W)=F_\ast+\frac12\tr\big((W-W_\ast)\Sigma(W-W_\ast)^\top\big)$ with $\Sigma\in\R^{d\times d}$ symmetric positive
definite: multi-output least squares, $\Sigma=\hat\Sigma+\lambda I$, $\hat\Sigma=\frac1n\sum_ix_ix_i^\top$; then
$G=E\Sigma$. (GLM) $F(W)=\frac1n\sum_i\ell_i(Wx_i)+\frac\lambda2\nF W^2$ with convex $C^2$ losses,
$c_i=\sup_z\lmax(\nabla^2\ell_i(z))<\infty$, and $f_i(W)=\ell_i(Wx_i)+\frac\lambda2\nF W^2$: per-sample loss gradients
have rank at most one.

\paragraph{Status.} \textbf{[P]} proved here (complete proof given). \textbf{[K]} known; the source is given and the
proof, when included, is for completeness. \textbf{[F]} formal computation (an expansion whose remainder is not
controlled here). \textbf{[N]} numerical: supported by the script named with the statement, not proved.
Table~\ref{tab:status} lists all statements with their status and the closest published statements. Statement numbers of
published papers are those read by the referees of 16 November 2026 in the versions they cite; they should be re-read
before being quoted elsewhere.

\section{Condition numbers in the spectral geometry}\label{sec:kappa}

In this section $H$ is a symmetric positive definite linear operator on $\R^{K\times d}$ (for the Frobenius inner
product), $q(\Delta)=\inner{\Delta}{H\Delta}$, and $\Lop(H)=\sup_{\Delta\ne0}q(\Delta)/\nop\Delta^2$,
$\muop(H)=\inf_{\Delta\ne0}q(\Delta)/\nop\Delta^2$. For a $C^2$ function, $\Lop(F)=\sup_W\Lop(\nabla^2F(W))$ and
$\muop(F)=\inf_W\muop(\nabla^2F(W))$. That $\Lop(H)$ is the smoothness constant for the spectral norm is known (for any
norm: Balles, Pedregosa and Le Roux, arXiv:2002.08056, Proposition~3; for the spectral norm: Shen et al.,
arXiv:2505.23737; Jawanpuria, Garai and Mishra, arXiv:2610.27747, Appendix~C, Remark~2; Islamov et al.,
arXiv:2603.05002, Eq.~(19)).

\begin{proposition}[variational formulas]\label{prop:var} \textbf{[P]}
\begin{enumerate}[label=(\alph*)]
\item $\displaystyle \Lop(H)=\max_{Q\in\cO}\inner{Q}{HQ}$.
\item $\displaystyle \frac1{\muop(H)}=\max_{\norm u=\norm v=1}\inner{uv^\top}{H^{-1}(uv^\top)}$; if $(u_\ast,v_\ast)$
attains the maximum, $\Delta_\ast=H^{-1}(u_\ast v_\ast^\top)$ attains $\muop(H)=q(\Delta_\ast)/\nop{\Delta_\ast}^2$.
\end{enumerate}
\end{proposition}
\begin{proof}
(a) Let $\nop\Delta\le1$, $\Delta=U\diag(s)V^\top$ with $s\in[0,1]^r$. The cube $[-1,1]^r$ is the convex hull of
$\{-1,1\}^r$, so $s=\sum_\varepsilon c_\varepsilon\varepsilon$ and $\Delta=\sum_\varepsilon c_\varepsilon U\diag(\varepsilon)V^\top$ is a convex
combination of elements of $\cO$. Since $q$ is convex, $q(\Delta)\le\max_{Q\in\cO}q(Q)$; $\cO$ is compact.

(b) We show that, for $\mu>0$, [$q(\Delta)\ge\mu\nop\Delta^2$ for all $\Delta$] is equivalent to
[$\inner{Y}{H^{-1}Y}\le\nnuc Y^2/\mu$ for all $Y$]. If the first holds, let $\Delta=H^{-1}Y$:
$\inner Y{H^{-1}Y}=\inner Y\Delta\le\nnuc Y\nop\Delta\le\nnuc Y\sqrt{q(\Delta)/\mu}$ and $q(\Delta)=\inner Y{H^{-1}Y}$.
If the second holds, take $Y=u_1v_1^\top$, a top singular pair of $\Delta$:
$\nop\Delta=\inner Y\Delta=\inner{H^{-1/2}Y}{H^{1/2}\Delta}\le\sqrt{\inner Y{H^{-1}Y}}\sqrt{q(\Delta)}\le\sqrt{q(\Delta)/\mu}$.
Hence $1/\muop=\sup_{\nnuc Y\le1}\inner Y{H^{-1}Y}$, the maximum of a convex function over the nuclear ball, attained at
an extreme point $\pm uv^\top$. At a maximizer $Y_\ast$, optimality gives $\inner{\Delta_\ast}{Y-Y_\ast}\le0$ on the ball
with $\Delta_\ast=H^{-1}Y_\ast$, i.e.\ $\inner{\Delta_\ast}{Y_\ast}=\nop{\Delta_\ast}$; then
$q(\Delta_\ast)=\inner{\Delta_\ast}{Y_\ast}=\nop{\Delta_\ast}=1/\muop$ and $q(\Delta_\ast)/\nop{\Delta_\ast}^2=\muop$.
\end{proof}

Item (b) is the classical duality between strong convexity for a norm and smoothness of the conjugate for the dual
norm (Kakade, Shalev-Shwartz and Tewari, JMLR 2012, Theorem~3) applied to a quadratic, with the reduction to extreme
points; its $\ell_1$ analogue is in Nutini et al.\ (ICML 2015, Section~4). The formulas are not a polynomial-time
recipe: maximizing a convex quadratic over the spectral ball contains the $\ell_\infty$ case (diagonal matrices), which
is NP-hard in general (Islamov et al., Appendix~A; Balles et al., Section~4.2). The values printed by
\texttt{check\_condition\_numbers.py} for operators without structure are therefore values of a multi-start ascent, i.e.\
lower bounds of $\Lop$ and of $1/\muop$; they are exact for the structured operators of Proposition~\ref{prop:struct}.

\begin{proposition}[frame bound, dimension barrier, rigidity]\label{prop:frame} \textbf{[P]; the barrier $\kop\ge r$ is [K]}
Let $(u_k)_{k\le r}$ be orthonormal in $\R^K$, $(v_k)_{k\le r}$ orthonormal in $\R^d$, and $h_k=q(u_kv_k^\top)$. Then
\[
\Lop(H)\ \ge\ \sum_{k\le r}h_k,\qquad \muop(H)\ \le\ \min_{k\le r}h_k,\qquad
\kop(H)\ \ge\ \frac{\sum_kh_k}{\min_kh_k}\ \ge\ r .
\]
Consequently $\kop(F)\ge r=\min(K,d)$ for every $C^2$ function satisfying \eqref{eq:consts}. Moreover $\kop(H)=r$ if
and only if $H$ is a multiple of the identity; hence $\kop(F)=r$ if and only if
$F=F_\ast+\frac\mu2\nF{W-W_\ast}^2$.
\end{proposition}
\begin{proof}
$\nop{u_kv_k^\top}=1$ gives $\muop\le h_k$. For signs $\varepsilon\in\{\pm1\}^r$ let
$Q_\varepsilon=\sum_k\varepsilon_ku_kv_k^\top\in\cO$. Averaging over uniform signs,
$\E\,q(Q_\varepsilon)=\sum_{k,l}\E[\varepsilon_k\varepsilon_l]\inner{u_kv_k^\top}{Hu_lv_l^\top}=\sum_kh_k$, so some $Q_\varepsilon$ has
$q(Q_\varepsilon)\ge\sum_kh_k$. For a function, apply this to $\nabla^2F(W)$ at any $W$.
\emph{Rigidity} (argument of the referee of Sections 2--3). If $\Lop=r\muop$, then for every frame
$r\muop\le\sum_kh_k\le\Lop$ forces $h_k=\muop$: $q(uv^\top)=\muop$ for all unit $u,v$. Moreover
$q(Q_\varepsilon)=r\muop+\sum_{k\ne l}\varepsilon_k\varepsilon_l\inner{u_kv_k^\top}{Hu_lv_l^\top}\le r\muop$ for all signs while the
cross terms have mean zero, so the cross terms vanish for all signs, i.e.\ $\inner{u_kv_k^\top}{Hu_lv_l^\top}=0$ for
$k\ne l$. With coordinate frames this gives $H_{ia,jb}=0$ for $i\ne j$, $a\ne b$; polarizing $q(uv^\top)=\muop\norm
u^2\norm v^2$ gives $H_{ia,ib}=\muop\delta_{ab}$ and $H_{ia,ja}=\muop\delta_{ij}$. So $H=\muop\,\mathrm{Id}$. For a
function, apply this at every $W$.
\end{proof}

The barrier itself is the matrix form of a folklore fact: a function that is $1$-strongly convex for $\ell_\infty$ on
$[-1,1]^n$ has range at least $n/2$ (Sidford and Tian, arXiv:1808.01278, Lemma~A.1), equivalently the cube is at
Banach--Mazur distance $\sqrt n$ from the Euclidean ball; restricting $F$ to $\{U\diag(s)V^\top\}$ reduces the spectral
case to it. The frame bound was not found; the rigidity was added during the revision and has not been searched for.
The frame bound can be loose: $q=\nF\Delta^2+2c\det\Delta$
on $2\times2$ matrices has all $h_k$ equal and $\kop=2/(1-\abs c)$.

\begin{corollary}[comparison with the Euclidean condition number]\label{cor:compare} \textbf{[P]}
For every eigenpair $(\lambda,V)$ of $H$: $\Lop(H)\ge\lambda\,\nr(V)$ and $\muop(H)\le\lambda\,\srank(V)$. With
$V_{\max},V_{\min}$ eigen-matrices for $\lmax(H)$ and $\lmin(H)$,
\[
\max\Big\{r,\ \frac{\kF(H)}{r},\ \kF(H)\,\frac{\nr(V_{\max})}{\srank(V_{\min})}\Big\}\ \le\ \kop(H)\ \le\ r\,\kF(H).
\]
In particular $\kop<\kF$ requires $\srank(V_{\min})>\nr(V_{\max})$: flat directions of higher rank than the sharp
ones. Both cases occur: $\kop\ge\kF$ whenever $V_{\min}$ can be chosen of rank one, and the lower bound $\kF/r$ is
approached (Example~\ref{ex:sharp}).
\end{corollary}
\begin{proof}
$\lambda\nnuc V=\inner{\polar(V)}{HV}\le\sqrt{q(\polar V)}\sqrt{q(V)}\le\sqrt{\Lop}\,\sqrt\lambda\,\nF V$ (Cauchy--Schwarz
for $q$, $\nop{\polar V}=1$), which is $\Lop\ge\lambda\nr(V)$ (sharpening due to the referee; the first version had
$\srank$). $\muop\le q(V)/\nop V^2=\lambda\srank(V)$. The other bounds are Proposition~\ref{prop:frame} and the comparison
of the norms.
\end{proof}

\begin{example}[the bound $\kF/r$ is sharp]\label{ex:sharp} \textbf{[P]} (Referee of Sections 2--3; it replaces a
numerical $\ell_\infty$ example of the first version whose printed bound was not established.) On $\R^{r\times r}$, for
$0<b\le m\le1$, let $q(\Delta)=m\nF\Delta^2+(1-m)\Delta_{12}^2-(m-b)(\tr\Delta)^2/r$: eigenvalues $1$ (on $E_{12}$),
$b$ (on $I/\sqrt r$) and $m$, so $\kF=1/b$. Then $\Lop=1+m(r-1)$ (for $Q$ orthogonal
$q(Q)=mr+(1-m)Q_{12}^2-(m-b)(\tr Q)^2/r\le mr+1-m$, with equality for a cyclic permutation with $Q_{12}=1$) and
$\muop=rmb/(m+(r-1)b)$ (attained at $\diag(1,c,\dots,c)$, $c=(m-b)/(m+(r-1)b)$). Hence
$\kop/\kF=(1+m(r-1))(m+(r-1)b)/(rm)$, which tends to $1/r$ when $b/m\to0$ and then $m\to0$; for $r=10$, $m=10^{-2}$,
$b=10^{-5}$, $\kop=\kF/9.09$. Checked in \texttt{check\_review\_additions.py}.
\end{example}

\begin{proposition}[structured Hessians]\label{prop:struct} \textbf{[P]; see Remark~\ref{rem:kappa} for what is [K]}
\begin{enumerate}[label=(\alph*)]
\item \emph{Kronecker.} If $H\Delta=B\Delta A$ with $A\in\R^{d\times d}$, $B\in\R^{K\times K}$ symmetric positive definite
with eigenvalues $a_1\ge\dots\ge a_d$, $b_1\ge\dots\ge b_K$, then
$\Lop=\sum_{k\le r}a_kb_k$, $\muop=a_db_K$, so that $\kop/\kF=\sum_{k\le r}a_kb_k/(a_1b_1)\in[1,r]$.
\item \emph{Least squares.} For (LS) with $\Sigma=\hat\Sigma+\lambda I$: $\muop(F)=\lmin(\Sigma)$ and $\Lop(F)$ is the sum
of the $r$ largest eigenvalues of $\Sigma$; with $f_i(W)=\frac12\norm{Wx_i-y_i}^2+\frac\lambda2\nF W^2$,
$\Lop(f_i)=\norm{x_i}^2+\lambda r$. If $K\ge d$,
\[
\Lop(F)=\tr\Sigma=\frac1n\sum_i\Lop(f_i),\quad\text{whereas}\quad
L_F(F)=\lmax(\hat\Sigma)+\lambda\le\frac1n\sum_iL_F(f_i)=\tr\hat\Sigma+\lambda :
\]
\emph{no averaging gain}, and $\kop(F)=\kF(F)\cdot\tr(\Sigma)/\lmax(\Sigma)$. If $K<d$ the averaging gain reappears
($\Lop(F)=\sum_{k\le K}\lambda_k(\Sigma)<\frac1n\sum_i\Lop(f_i)$ unless $\hat\Sigma$ has rank at most $K$), and for $K=1$ the
spectral norm is the Euclidean norm.
\item \emph{Ridge-regularized models: no free lunch, for every norm.} For (GLM), assume there is a sequence $W_t$ along
which $\max_i\nop{\nabla^2\ell_i(W_tx_i)}\to0$ (softmax and logistic losses: $W_t=tW_1$ for a generic $W_1$). Let $N$ be
\emph{any} norm on $\R^{K\times d}$ and $c_N=\max_{\nF\Delta=1}N(\Delta)$. Then, for every convex $C^2$ function $g$,
\[
\mu_N(F)=\frac{\lambda}{c_N^2}=\frac{\mu_F(F)}{c_N^2},\qquad L_N(g)\ \ge\ \frac{L_F(g)}{c_N^2},
\]
hence
\[
\frac{L_N(F)}{\mu_N(F)}\ \ge\ \kF(F),\qquad \frac{L_N(f_i)}{\mu_N(F)}\ \ge\ \frac{L_F(f_i)}{\mu_F(F)}=\frac{c_i\norm{x_i}^2+\lambda}{\lambda}\quad\text{for every } i :
\]
no norm has a better worst-case condition number than the Frobenius norm, for the full objective, and for the maximum
or the mean of the per-sample constants that govern variance-reduced methods.
\item \emph{Instances of (c).} Spectral norm ($c_N=1$): $\muop(F)=\mu_F(F)=\lambda$ and $\Lop(f_i)=c_i\norm{x_i}^2+\lambda
r$. Right- and two-sided preconditioned spectral norms $N(\Delta)=\nop{S^{1/2}\Delta P^{1/2}}$ ($c_N^2=\lmax(S)\lmax(P)$),
e.g.\ $\mu_{N_P}(F)=\lambda/\lmax(P)$ and, with $y_i=P^{-1/2}x_i$,
$L_{N_P}(f_i)=\sum_{k\le r}\lambda_k\big(c_iy_iy_i^\top+\lambda P^{-1}\big)$ for $N_P(\Delta)=\nop{\Delta P^{1/2}}$.
Euclidean metrics $\nF{\Delta P^{1/2}}$ and their Kronecker-factored versions. For $K=1$ and the logistic loss, the
last case is Proposition B1-nofreelunch of the note B1.
\end{enumerate}
\end{proposition}
\begin{proof}
(a) For $Q\in\cO$, $q(Q)=\tr(A\,Q^\top BQ)\le\sum_k\lambda_k(A)\lambda_k(Q^\top BQ)$ (von Neumann's trace inequality).
If $K\ge d$, $Q^\top Q=I_d$ and $\lambda_k(Q^\top BQ)\le b_k$ (Poincar\'e separation); if $K<d$, $QQ^\top=I_K$ and the
nonzero eigenvalues of $Q^\top BQ$ are those of $B$. In both cases $q(Q)\le\sum_{k\le r}a_kb_k$, with equality for
$Q=\sum_{k\le r}u^B_k(u^A_k)^\top$. Next, $q(\Delta)=\nF{B^{1/2}\Delta A^{1/2}}^2\ge b_Ka_d\nF\Delta^2\ge b_Ka_d\nop\Delta^2$,
with equality for $\Delta=u^B_K(u^A_d)^\top$.
(b) is (a) with $A=\Sigma$, $B=I$; $D^2f_i[\Delta,\Delta]=\norm{\Delta x_i}^2+\lambda\nF\Delta^2$ is at most
$\norm{x_i}^2+\lambda r$ on the spectral unit ball, with equality for an element of $\cO$ whose row space contains $x_i$,
and at most $(\norm{x_i}^2+\lambda)\nF\Delta^2$, with equality for $\Delta=ux_i^\top/\norm{x_i}$.
(c) For every $W$ and $\Delta$, $D^2F(W)[\Delta,\Delta]\ge\lambda\nF\Delta^2\ge(\lambda/c_N^2)N(\Delta)^2$; at a maximizer
$\Delta_0$ of $N/\nF\cdot$, $D^2F(W_t)[\Delta_0,\Delta_0]\le\max_i\nop{\nabla^2\ell_i(W_tx_i)}\frac1n\sum_i\norm{\Delta_0x_i}^2
+\lambda\nF{\Delta_0}^2\to\lambda\nF{\Delta_0}^2$. So $\mu_N(F)=\lambda/c_N^2$. For a convex $g$,
$D^2g[\Delta,\Delta]/N(\Delta)^2\ge D^2g[\Delta,\Delta]/(c_N^2\nF\Delta^2)$; take suprema. $L_F(f_i)=c_i\norm{x_i}^2+\lambda$
is attained in the direction $ex_i^\top/\norm{x_i}$ with $e$ a top eigenvector of $\nabla^2\ell_i$.
(d) $\Lop(f_i)$: an element of $\cO$ mapping $x_i/\norm{x_i}$ to a top eigenvector of $\nabla^2\ell_i$. For $N_P$, by
Proposition~\ref{prop:var}(a) one maximizes over $\Delta'=\Delta P^{1/2}\in\cO$; for $K\le d$ with row space $S$,
$(\Delta'y_i)^\top M(\Delta'y_i)+\lambda\tr(\Delta'P^{-1}\Delta'^\top)\le\tr\big((\lmax(M)y_iy_i^\top+\lambda P^{-1})\Pi_S\big)$
with equality when $\Delta'$ maps $\Pi_Sy_i$ to a top eigenvector of $M=\nabla^2\ell_i$; maximize over $S$ (Ky Fan); for
$K\ge d$ the full trace is obtained.
\end{proof}

\begin{remark}[reading and attribution]\label{rem:kappa}
(i) \emph{What is known.} For a least-squares layer, $\Lop=\tr\Sigma$, $L_F=\lmax(\Sigma)$ and the reading of their ratio
as the stable rank of the data are in Davis and Drusvyatskiy (arXiv:2512.04299, Eq.~(1.5) and Section~1.2.1); the sum of
the $K$ largest eigenvalues for $K<d$ is in Jawanpuria et al.\ (Appendix~C, Remark~2); the bound $\sum_ka_kb_k$ for a
Kronecker Hessian, by von Neumann's inequality, is in Shen et al.\ (arXiv:2505.23737, Section~4.3.1), and the trace norm
as non-Euclidean smoothness constant in Kovalev and Borodich (arXiv:2511.11466, Lemma~2). The rate
$1-\lmin(H)/\tr(H)$ of the nuclear-scaled step on matrix quadratics is in Jawanpuria et al.\ (Theorem~3) and Shulgin,
Laing, Orvieto and Richt\'arik (arXiv:2610.30280, Theorem~3, Eq.~(16), tight); PolarGrad has the weaker factor $r\kappa$
with Frobenius constants (Theorem~3.2 of v4). For $\ell_\infty$, comparisons of $L_\infty$ with $L_2$ are in Balles et
al.\ (Propositions 4--5, Eq.~(22)) and of $\mu_1$ with $\mu_2$ in Nutini et al. Not found: the frame bound,
Corollary~\ref{cor:compare} in its first form (with the stable rank on both sides), the attainment and $\muop$ in (a),
the identity ``$\Lop(F)$ = mean of the $\Lop(f_i)$'' read as an absence of averaging gain, and items (c)--(d) for the
spectral norm and its preconditioned versions. Not searched (added during the revision): the rigidity, the nuclear-rank
form of Corollary~\ref{cor:compare}, Example~\ref{ex:sharp}, and item (c) for an arbitrary norm.

(ii) \emph{No averaging gain.} In the Euclidean geometry the smoothness of an average is smaller than the average
smoothness because the curvatures $x_ix_i^\top$ of different samples point in different directions; an isometry
$\Delta$ excites all samples at once, so for $K\ge d$ the triangle inequality is an equality. Then
$\kop(F)=\bar L_{\mathrm{op}}/\mu$ with $\bar L_{\mathrm{op}}=\frac1n\sum_i\Lop(f_i)=\tr\hat\Sigma+\lambda d$: this is the
constant $\bar L_F/\mu$ of variance-reduced methods with importance sampling ($\bar L_F=\tr\hat\Sigma+\lambda$; Xiao and
Zhang, SIAM J.\ Optim.\ 2014, arXiv:1403.4699, Theorem~1 and the remarks that follow it) plus
$(d-1)\lambda/\mu$, and equal to it when $\lambda=0$. The full-batch nuclear-scaled step has, up to that term, the
condition number of an importance-sampled variance-reduced method, but pays $n$ gradients per step.

(iii) \emph{Counterpart of B1 and B3.} Item (c) says nothing specific about the spectral geometry: it is one
norm-independent fact, with a three-line proof, of which Proposition B1-nofreelunch (Euclidean metrics $P$, $K=1$) is
an instance; the referee of Sections 2--3 and the author arrived at it independently during the revision. The mechanism
is that the curvature floor of a ridge-regularized model is isotropic. For softmax the hypothesis is not even needed for
the spectral norm: $\nabla^2\ell_i(z)\ones=0$, so the gauge direction $\ones v^\top$ has curvature exactly
$\lambda\nF\cdot^2$ at every $W$. Any advantage of a polar step on such problems must therefore be a property of the
trajectory, as for preconditioned SAGA in B1. Sections~\ref{sec:dual}--\ref{sec:ls} make this quantitative.
\end{remark}

\section{Polar steps as mirror descent on the nuclear norm of the gradient}\label{sec:dual}

Assume \eqref{eq:consts} with $\muop>0$. Then $\nabla F$ is a bijection of $\R^{K\times d}$, with inverse
$\nabla F^\ast$, where $F^\ast$ is the convex conjugate of $F$; $F^\ast$ is $\frac1{\Lop}$-strongly convex and
$\frac1{\muop}$-smooth for the nuclear norm (Kakade et al., Theorem~3). Write
$D_\Phi(a,b)=\Phi(a)-\Phi(b)-\inner{\nabla\Phi(b)}{a-b}$.

\begin{proposition}[dual representation]\label{prop:dual} \textbf{[K], except the nonsmooth reading of (a)}
\begin{enumerate}[label=(\alph*)]
\item With $G=\nabla F(W)$, the polar step $W^+=W-\eta O$, $O\in\partial\psi(G)$, is
\[
\nabla F^\ast(G^+)=\nabla F^\ast(G)-\eta\,O,\qquad O\in\partial\psi(G):
\]
one step of mirror \emph{subgradient} descent on the convex, $1$-Lipschitz function $\psi=\nnuc\cdot$, in the space of
gradients, with mirror map $F^\ast$. The minimizer of $\psi$ is $G=0$, i.e.\ $W=W_\ast$, and $\min\psi=0$ is known.
\item $D_{F^\ast}(G',G)=D_F(W,W')$ for $G=\nabla F(W)$, $G'=\nabla F(W')$; in particular $D_{F^\ast}(0,G)=F(W)-F_\ast$.
\item (Dual sharpness and smoothness.) $2\muop\,(F(W)-F_\ast)\le\nnuc{\nabla F(W)}^2\le2\Lop\,(F(W)-F_\ast)$.
\item (One step.) $F(W^+)-F_\ast\le F(W)-F_\ast-\eta\nnuc G+\frac{\Lop}2\eta^2$.
\end{enumerate}
\end{proposition}
\begin{proof}
(a) $W=\nabla F^\ast(G)$ and $W^+=\nabla F^\ast(G^+)$. (b) $F^\ast(G)=\inner WG-F(W)$; substitute. (c) Strong convexity
gives, for all $W'$, $F(W')\ge F(W)-\nnuc G\nop{W'-W}+\frac{\muop}2\nop{W'-W}^2\ge F(W)-\nnuc G^2/(2\muop)$; take
$W'=W_\ast$. Smoothness gives $F_\ast\le F(W-t\,\polar(G))\le F(W)-t\nnuc G+\frac{\Lop}2t^2$; take $t=\nnuc G/\Lop$. (d) is
smoothness with $\inner GO=\nnuc G$, $\nop O\le1$.
\end{proof}

For a quadratic $F$ with Hessian $H$ and $G\ne0$, write (the quantity $J$ is the $J_t$ of Shen et al., Theorem~4.14)
\[
J(W)=\inner{O}{HO}\in[\muop,\Lop]\quad(O=\polar(G)),\qquad
\mu_{\mathrm{PL}}(W)=\frac{\nnuc G^2}{2(F(W)-F_\ast)}\in[\muop,J(W)].
\]

\begin{proposition}[exact one-step representation]\label{prop:repr} \textbf{[K] for (LS); immediate in general}
For a quadratic $F$ and the step $W^+=W-\frac{\nnuc G}{L}\polar(G)$,
\begin{equation}\label{eq:repr}
\frac{F(W^+)-F_\ast}{F(W)-F_\ast}=1-\frac{\mu_{\mathrm{PL}}(W)}{L}\Big(2-\frac{J(W)}{L}\Big),
\end{equation}
and the best $L$ is $J(W)$ (exact line search), with ratio $1-\mu_{\mathrm{PL}}(W)/J(W)\ge0$. With $L=\Lop$, the
certified factor $1-1/\kop$ is the value of \eqref{eq:repr} at the worst cases $J=\Lop$, $\mu_{\mathrm{PL}}=\muop$. The
slack of the smoothness step can improve the exponent by a factor at most $2$ (since $2-J/\Lop\in[1,2)$); the slack of
the sharpness step, $\mu_{\mathrm{PL}}/\muop\in[1,\kop]$, is the only unbounded one. Moreover
$\mu_{\mathrm{PL}}(W)=\nr(G)\cdot\nF G^2/(2(F(W)-F_\ast))$ and
$\inner GE/\nop E^2=\srank(E)\cdot\inner GE/\nF E^2$: the spectral effective constants are the Euclidean ones of the
note B2 multiplied by the nuclear rank of the gradient and by the stable rank of the error.
\end{proposition}
\begin{proof}
$F(W-\eta O)=F(W)-\eta\inner GO+\frac{\eta^2}2\inner O{HO}$ exactly; substitute $\eta=\nnuc G/L$ and divide by
$F(W)-F_\ast=\nnuc G^2/(2\mu_{\mathrm{PL}})$.
\end{proof}

For the one-sided quadratic, \eqref{eq:repr} is Eq.~(17) of Theorem~3 of Shulgin et al.\ (with $\alpha=J/L$,
$p=\mu_{\mathrm{PL}}/J$) and Eq.~(8) of Jawanpuria et al. Its interest here is the parallel with SAGA in the note B3
(Proposition B3-representation): the proof of the certified rate is a worst-case substitution into an exact formula,
the gain available from the smoothness inequality is bounded (there $8/3$, here $2$), and the curvature met along the
trajectory is the unbounded factor.

\begin{theorem}[step rules]\label{thm:steps} \textbf{[P]; the rate of (a) is [K]}
Let $D_t=F(W_t)-F_\ast$ along polar steps $W_{t+1}=W_t-\eta_tO_t$, $O_t\in\partial\psi(G_t)$.
\begin{enumerate}[label=(\alph*)]
\item \emph{Polyak step = nuclear-scaled step.} The mirror Polyak step for $\psi$ (optimal value $0$, mirror map $F^\ast$
of modulus $1/\Lop$, subgradients of dual norm $1$) is $\eta_t=\psi(G_t)/\Lop$; it gives $D_{t+1}\le(1-1/\kop)D_t$.
\item \emph{Geometric schedule.} Let $\bar D_0\ge D_0$, $\rho=\max\{1/\kop,\,1-1/\kop\}$ (equal to $1-1/\kop$ as soon as
$r\ge2$, by Proposition~\ref{prop:frame}) and $\eta_t=\eta_0\rho^{t/2}$ with $\eta_0=\sqrt{2\muop\bar D_0}/\Lop$. Then
$D_t\le\bar D_0\rho^t$ for all $t$: a fixed-length schedule reaches the rate of the nuclear-scaled step, at the price of
knowing $\muop$ and a bound on $D_0$.
\item \emph{Constant length.} If $\eta_t=\eta$, then for every $t$ with $G_t\ne0$, $\max(D_t,D_{t+1})\ge\muop\eta^2/8$
(the constant is attained: $F=\frac\mu2w^2$, $w_0=\eta/2$), and
$\limsup_tD_t\le\max\{\kop/8,1/2\}\,\Lop\eta^2$.
\end{enumerate}
\end{theorem}
\begin{proof}
By Proposition~\ref{prop:dual}(c)--(d), $D_{t+1}\le\varphi_{\eta_t}(D_t)$ with
$\varphi_\eta(x)=x-\eta\sqrt{2\muop x}+\frac{\Lop}2\eta^2$, a convex function of $x\ge0$.
(a) Minimize $-\eta\nnuc{G_t}+\frac\Lop2\eta^2$ directly:
$D_{t+1}\le D_t-\nnuc{G_t}^2/(2\Lop)\le(1-\muop/\Lop)D_t$.
(b) Induction: if $D_t\le\bar D_t:=\bar D_0\rho^t$ and $\eta_t=\sqrt{2\muop\bar D_t}/\Lop$, convexity gives
$D_{t+1}\le\max\{\varphi_{\eta_t}(0),\varphi_{\eta_t}(\bar D_t)\}=\max\{\bar D_t/\kop,\ \bar D_t(1-1/\kop)\}=\rho\bar D_t$;
and $\eta_{t+1}=\eta_t\sqrt\rho$.
(c) If $G_t\ne0$, $\eta=\nop{W_{t+1}-W_t}\le\nop{E_t}+\nop{E_{t+1}}$ and $D\ge\frac{\muop}2\nop E^2$ give the lower bound.
For the upper bound, $\varphi_\eta(x)\le x$ iff $x\ge x_0:=\Lop^2\eta^2/(8\muop)$, and
$x_1:=\max\{\varphi_\eta(0),\varphi_\eta(x_0)\}=\max\{\Lop\eta^2/2,x_0\}$. The interval $[0,x_1]$ is invariant: if
$D_t\le x_0$ then $D_{t+1}\le x_1$ by convexity, and if $x_0<D_t\le x_1$ then $D_{t+1}\le\varphi_\eta(D_t)\le D_t$. If
$D_t>x_1$ for all $t$, then $(D_t)$ is nonincreasing and converges to some $D_\infty\ge x_1\ge x_0$; passing to the
limit in $D_{t+1}\le\varphi_\eta(D_t)$ gives $D_\infty\le\varphi_\eta(D_\infty)$, i.e.\ $D_\infty\le x_0$, so
$D_\infty=x_0=x_1$. In both cases $\limsup D_t\le x_1$.
\end{proof}

\begin{proposition}[flow]\label{prop:flow} \textbf{[K] in substance; the spectral instance is written out here}
Let $W(\cdot)$ be an absolutely continuous solution of the differential inclusion
$\dot W(t)\in-\partial\psi(\nabla F(W(t)))$. Then, for almost every $t$ with $G(t)\ne0$,
\[
\frac{d}{dt}F(W(t))=-\nnuc{G(t)},\qquad \frac{d}{dt}\nnuc{G(t)}=-J(t),\quad
J(t)=D^2F(W(t))[O_t,O_t]\in[\muop,\Lop],
\]
where $O_t=-\dot W(t)$. Hence $G$ vanishes at a finite time $T$ (and $W(t)=W_\ast$ for $t\ge T$) with
\begin{equation}\label{eq:sandwich}
\nop{W_0-W_\ast}\ \le\ T\ \le\ \sqrt{\frac{2(F(W_0)-F_\ast)}{\muop}}\ \le\ \sqrt{\kop}\,\nop{W_0-W_\ast}.
\end{equation}
For a quadratic $F$ the solution exists and is unique (in the variable $G$ it is the gradient flow of $\psi$ for the
inner product $\inner\cdot{H^{-1}\cdot}$), $t\mapsto\nnuc{G(t)}$ is convex and $J(t)$ is nonincreasing.
\end{proposition}
\begin{proof}
$\frac d{dt}F(W)=\inner G{\dot W}=-\inner G{O_t}=-\nnuc G$. The curves $G(\cdot)$ and $\psi(G(\cdot))$ are Lipschitz on
bounded intervals, hence differentiable at almost every $t$; at such a $t$ and for $Z\in\partial\psi(G(t))$,
$\psi(G(t+h))-\psi(G(t))\ge\inner Z{G(t+h)-G(t)}$ for $h$ of both signs, so $\frac d{dt}\psi(G(t))=\inner Z{\dot G(t)}$;
take $Z=O_t$ and $\dot G=-\nabla^2F(W)[O_t]$. Since $\nop{O_t}=1$ when $G\ne0$, the bounds on $J$ follow from
\eqref{eq:consts}. With $D=F(W)-F_\ast$ and Proposition~\ref{prop:dual}(c), $\dot D=-\nnuc G\le-\sqrt{2\muop D}$, i.e.\
$\frac d{dt}\sqrt D\le-\sqrt{\muop/2}$: extinction before $\sqrt{2D_0/\muop}$, and $D_0\le\frac\Lop2\nop{E_0}^2$. After
$T$, $F=F_\ast$ since $F$ is nonincreasing. Finally $\nop{\dot W}\le1$ gives $\nop{W_0-W_\ast}\le T$. For a quadratic,
$\dot G\in-H\partial\psi(G)$ is the gradient flow of a convex function in a Hilbert space: existence, uniqueness,
convexity of $t\mapsto\psi(G(t))$ and monotonicity of the speed $\inner{O_t}{HO_t}^{1/2}$ are Theorems 3.1--3.2 of
Br\'ezis (\emph{Op\'erateurs maximaux monotones}, 1973), as quoted by Bungert, Burger, Chambolle and Novaga
(Analysis \& PDE 2021, Theorems 4.1--4.2).
\end{proof}

\begin{remark}[attribution and use]\label{rem:dual}
\emph{Dual reading.} Methods $W^+=W-\eta\nabla k(\nabla F(W))$ are the dual-space preconditioned gradient methods of
Maddison, Paulin, Teh and Doucet (SIAM J.\ Optim.\ 2021). Their Section~5 already describes the method as ``a Bregman
gradient method in the dual space, which minimizes the designed $k$ while the conjugate $f^\ast$ plays the role of the
reference function'', names steepest descent for a smooth norm as the case $k=\frac12\norm\cdot_\ast^2$, and states that
the analysis ``does not apply in the case of other norms or normalized steepest descent''; item (b) is their Lemma~3.2.
The same method is the dual mirror descent of Kim et al.\ (arXiv:2311.17296). Chen, Li and Liu (arXiv:2506.15054) identify
Muon's direction with a subgradient of the nuclear norm (Lion-$\mathcal K$), and Ghane, Akhtiamov and Hassibi
(arXiv:2603.10485) list nonsmooth preconditioners ``such as Muon'' as future work. What is added here is that nonsmooth
case carried out: $k=\psi$ (mirror subgradient descent on a sharp function) and $k=\frac12\psi^2$ (nuclear-scaled step).
\emph{Step rules.} The rate of (a) is known: Balles et al.\ (Proposition~9, any norm), Gluon (Riabinin et al.,
arXiv:2505.13416, Theorem~5, spectral constants), Nutini et al.\ and Karimi, Nutini and Schmidt for $\ell_1$/sign steps.
The mirror Polyak step is defined by D'Orazio et al.\ (arXiv:2110.15412); the identification with the nuclear-scaled step
was not found. Statement (b) is the geometric step-size rule of subgradient methods on sharp functions (Goffin, Math.\
Programming 1977; Shor) transported by Proposition~\ref{prop:dual}; the closest published statement is Theorem~4.1 of Yang,
He and Li (arXiv:2605.26977: spectral descent with a geometric schedule on convex, Lipschitz, sharp objectives, same
induction); geometric schedules for specific models are in Ma, Huang, Chi and Chen (arXiv:2601.13474) and Rhee, Simon and
Karkada (arXiv:2606.30509). For a general smooth strongly convex $F$ the statement was not found. Floors of
constant-length steps are in Jawanpuria et al.\ (Lemma~4), Gonon, Mu\c{s}at and Boumal (arXiv:2602.11948, Section~3.1),
PolarGrad (Theorem~3.4, an upper-bound recursion) and Vaswani, Sun and Babanezhad (arXiv:2605.30648, Proposition~8, the
two-cycle of value $\eta^2/8$); the explicit $\limsup$ bound was not found. \emph{Flow.} The inclusion for an arbitrary
norm, existence, the identity $\frac d{dt}F=-\norm{\nabla F}_\ast$ and the bound $T\le\sqrt{2D_0/\mu}$ under a dual-norm
Polyak--\L{}ojasiewicz inequality are in Leplat, Mayorga, Hildebrand and Gasnikov (arXiv:2508.18510, Theorem~3.1 and
Eq.~(16)--(18)); for sign flows see also Ma, Wu and E (MSML 2021) and Cort\'es (Automatica 2006), whose Theorem~9
(numbering of the preprint read by the referees) uses
$\norm{\nabla f}_1$ as Lyapunov function with derivative $-v^\top\nabla^2f\,v$: the $\ell_\infty$ form of the second
identity. For quadratics the extinction-time bounds are those of Bungert et al.\ for gradient flows of one-homogeneous
functionals (Theorem~6.2, Proposition~6.9). A continuous-time analysis of \emph{smoothed} polar flows is in Liu et al.\
(arXiv:2608.01911); it leaves the exact flow through rank changes for future work. The first version of these notes
stated the weaker bounds $\nnuc{G_0}/\Lop\le T\le\nnuc{G_0}/\muop$ and presented their ratio $\kop$ as intrinsic; both
are implied by \eqref{eq:sandwich}, whose ratio is $\sqrt{\kop}$. Each of the two weaker bounds is attained on (LS)
($E_0=au\,v_d^\top$ and $E_0=aQ$, $Q^\top Q=I_d$). The numerical checks are in \texttt{check\_dual\_mirror.py} and
\texttt{check\_review\_additions.py}.
\end{remark}

\section{The least-squares model: published dynamics and additions}\label{sec:ls}

Throughout this section $F$ is the one-sided quadratic (LS), $K\ge d$, $\Sigma$ has eigenvalues
$\lambda_1\ge\dots\ge\lambda_d>0$ with orthonormal eigenvectors $v_1,\dots,v_d$, $p_k=\lambda_k/\tr\Sigma$, and
$G=E\Sigma$. By Proposition~\ref{prop:struct}, $\Lop=\tr\Sigma$, $\muop=\lambda_d$, $\kop=1/p_d$. The first version of
these notes presented the dynamics of this model as new. The literature review showed that most of it is published,
in particular in the study of matrix quadratics by Shulgin, Laing, Orvieto and Richt\'arik (arXiv:2610.30280, 21 October
2026; their matrices are the transposes of ours, $H=\Sigma$ acting on the row side). Section~\ref{sec:ls-known} lists
what is known in our notation; Section~\ref{sec:ls-new} contains the additions.

\paragraph{Convention.} In this section a polar step uses $O_t=\polar(G_t)$, the thin polar factor, which vanishes on
the kernel of $G_t$: this is what makes the rank of the error non-increasing and gives $\sign(0)=0$ in (K3). For this
model it is also the direction selected by the flow. Indeed, for $Z=UV^\top+P_{U^\perp}CP_{V^\perp}\in\partial\psi(G)$,
\[
\tr(Z\Sigma Z^\top)=\tr(V^\top\Sigma V)+\nF{P_{U^\perp}CP_{V^\perp}\Sigma^{1/2}}^2
\]
(the cross term $2\tr(UV^\top\Sigma P_{V^\perp}C^\top P_{U^\perp})$ vanishes because $P_{U^\perp}U=0$), so $\polar(G)\Sigma$
is the element of least norm of $\partial\psi(G)\Sigma$ for the inner product $\inner{\cdot}{\cdot\,\Sigma^{-1}}$ of
Proposition~\ref{prop:flow}, hence minus the right derivative of $G(t)$ at \emph{every} time (same theorem of Br\'ezis).
With another element of the subdifferential (what a full singular value decomposition returns at an exactly singular
matrix) the zero amplitudes are re-excited at the scale of the step. This selection property, pointed out by the
mathematical referee of this section, is specific to one-sided Hessians: it fails for $\Delta\mapsto B\Delta A$.

\subsection{What is published}\label{sec:ls-known}

\begin{enumerate}[label=(K\arabic*)]
\item \emph{Constant curvature along the polar factor.} When $G$ has rank $d$, $O=\polar(G)$ satisfies $O^\top O=I_d$, so
\begin{equation}\label{eq:Jconst}
J(W)=\tr(O\Sigma O^\top)=\tr\Sigma=\Lop :
\end{equation}
the smoothness inequality of the spectral geometry is an equality at every full-rank iterate, the nuclear-scaled step
with $L=\tr\Sigma$ is the exact line search along the polar factor, and its contraction factor is $1-p(W)$ with
$p(W)=\nnuc G^2/(2\tr\Sigma\,(F-F_\ast))=(\tr\abs G)^2/(\tr\Sigma\,\tr(\Sigma^{-1}\abs G^2))\ge\lambda_d/\tr\Sigma$, the lower
bound being tight as a supremum, on errors that are nearly of rank one (Shulgin et al., Theorem~3, Eq.~(14)--(16) and
Appendix~E.2, Eq.~(17) for a general length, Corollary~3; Jawanpuria et al., Remark~2, Eq.~(8), Theorem~3; Wang, Zhang,
Li, Bergemann and Yang, arXiv:2606.04662, Eq.~(E.1); Davis and Drusvyatskiy, Eq.~(1.5)). At an iterate of rank $\rho<d$,
$J=\tr(\Sigma P)<\tr\Sigma$ with $P$ the projector on the row space of $G$ (Shulgin et al., Eq.~(18)), and for a
gradient of rank one the polar factor is proportional to the gradient: the two line searches coincide (their Remark~1).
\item \emph{One step along the polar factor.} With $G_0=O_0Y_0$, $Y_0=\abs{G_0}\succ0$: along $W_0-\gamma O_0$ the gradient
is $O_0(Y_0-\gamma\Sigma)$, $F=F(W_0)-\gamma\tr Y_0+\frac12\gamma^2\tr\Sigma$, the rank drops iff $\gamma$ is an eigenvalue of
$\Sigma^{-1/2}Y_0\Sigma^{-1/2}$, and the next root is $\abs{Y_0-\gamma\Sigma}$ (Shulgin et al., Propositions~9, 10, 12
and 17). The line-search step is exact in one step iff $E_0^\top E_0$ is a multiple of $I_d$ (their Theorem~1(iii)), and
after a line-search step $\tr(Y_0-\gamma\Sigma)=0$ (proof of their Proposition~18). For $d=2$ the line-search iteration
is solved completely from every full-rank start: after one step the root is a multiple of the identity and then $F$
contracts by $((\kappa-1)/(\kappa+1))^2$ per step (their Proposition~18, Corollary~6 and Remark~7).
\item \emph{Commuting case.} If $E_0^\top E_0$ commutes with $\Sigma$, write $E_0=\sum_ka_ku_kv_k^\top$ with $a_k\ge0$.
Then for any lengths $E_t=\sum_k\sigma_k(t)u_kv_k^\top$ with $\sigma_k(t+1)=\sigma_k(t)-\eta_t\sign(\sigma_k(t))$: the
iterates do not depend on $\Sigma$. This is Eq.~(4) of Gonon, Mu\c{s}at and Boumal for $\Sigma=I$; Proposition~1 and
Appendix~B.2 of Vasudeva, Deora, Zhao, Sharan and Thrampoulidis (arXiv:2510.22980) for this model from $W_0=0$ under joint
diagonalizability, including the flow $E(t)=\sum_k(a_k-t)_+u_kv_k^\top$ (singular-value soft-thresholding of the
initial error); Ma et al.\ for their models; Corollary~4 of Shulgin et al.\ (``$H$ drops out of the direction''; the
paragraph that follows it treats the fixed dual-adapted steps $s_i^+=\abs{s_i-\alpha\bar s_H}$ and notes that
commutation persists along the trajectory). The iterates
are independent of $\Sigma$ for lengths fixed in advance; the nuclear-scaled length depends on $\Sigma$ through $p$. With
the nuclear-scaled step the amplitudes follow $a\mapsto\abs{a-\E_p[a]}$, the mean being taken under the weights $p$,
and the one-step factor is $1-(\E_pa)^2/\E_p[a^2]$ (Shulgin et al., Eq.~(49)--(53) and Remark~4; Wang et al.,
Proposition~F.2), and the two-valued configurations $a=\delta((1-\pi)\ones_A+\pi\ones_{A^c})$, $\pi=p(A)$, are invariant
with $F$ contracting by $(1-2\pi)^2$ per step (Shulgin et al., Corollary~5; Wang et al., Proposition~E.3).
Rank-deficient invariant directions, on which the error is multiplied by $1-p(S)$ at each step ($F$ by $(1-p(S))^2$),
are in Jawanpuria et al.\ (Proposition~8). Related exact recursions in neighbouring models: Du and Su
(arXiv:2604.01472, Lemma~1 and Corollary~2) for a least-squares layer with a single-spike covariance (sign recursions
on the commuting modes, a normalized-gradient recursion on the mixed mode); Kang, Yun, Shin, Cho and Yun
(arXiv:2602.06385, Theorem~5.7 and Remark~6.8) for factorization models (equal-rate spectral dynamics, finite-time
convergence of the exact spectral gradient flow).
\item \emph{Slow regime and its remedy.} The line search subtracts a mean lying in the bulk, sends the amplitudes
near it close to zero and so creates the tiny singular values that make $p(W)$ small: ``a slow [regime], caused by tiny
singular values of the iterate, in which PolarGD is much slower than gradient descent'' (Shulgin et al., abstract and
Section~4.1). A finite number of Newton--Schulz iterations acts like gradient descent on singular values below
$1/\ell_5\approx2\cdot10^{-3}$ (their Section~4.2 and Appendix~G; Jawanpuria et al., Proposition~2, for the local rate of
bounded-gain maps with a constant step).
\end{enumerate}

\subsection{Additions}\label{sec:ls-new}

\begin{proposition}[first phase, for every schedule]\label{prop:phase1} \textbf{[P]; ingredients are (K2)}
Let $G_0=E_0\Sigma$ have rank $d$, $G_0=O_0Y_0$, and
\[
\tau_\ast=\lmin\big(\Sigma^{-1/2}Y_0\Sigma^{-1/2}\big),\qquad \tau_{\mathrm{ls}}=\frac{\tr Y_0}{\tr\Sigma}=\frac{\nnuc{G_0}}{\Lop}.
\]
\begin{enumerate}[label=(\alph*)]
\item For any polar steps $W_{t+1}=W_t-\eta_t\polar(G_t)$ with lengths $\eta_s\ge0$ and $\tau_t:=\sum_{s<t}\eta_s<\tau_\ast$,
the polar factor is frozen: $\polar(G_t)=O_0$ and $W_t=W_0-\tau_tO_0$. The same holds for a momentum with non-negative
weights, and the flow is the segment $W_0-tO_0$ on $[0,\tau_\ast]$: its full-rank phase lasts exactly $\tau_\ast$.
\item $s_d(E_0)\le\tau_\ast\le\min\big\{\big(\prod_is_i(E_0)\big)^{1/d},\ \tau_{\mathrm{ls}}\big\}\le\nop{E_0}$, and
$\tau_{\mathrm{ls}}=\tau_\ast$ iff $E_0^\top E_0$ is a multiple of $I_d$: the exact line search always leaves the first
phase unless it is exact in one step.
\item Afterwards $G_t=O_0M_t$ with $M_{t+1}=M_t-\eta_t\polar(M_t)\Sigma$, and after one line-search step
$M_1=Y_0-\tau_{\mathrm{ls}}\Sigma$ is symmetric with trace zero. Consequently the \emph{second} nuclear-scaled step
($L=\tr\Sigma$) from any full-rank start with $W_1\ne W_\ast$, commuting or not (also when $M_1$ is singular), satisfies
\[
\frac{F(W_2)-F_\ast}{F(W_1)-F_\ast}\ \le\ 1-\frac{4p_dp_{d-1}}{p_d+p_{d-1}}\ \le\ 1-\frac2{\kop}.
\]
\end{enumerate}
\end{proposition}
\begin{proof}
(a) If $W=W_0-\tau O_0$ then $G=O_0(Y_0-\tau\Sigma)$ with $Y_0-\tau\Sigma\succ0$ iff $\tau<\tau_\ast$; a product of a matrix
with orthonormal columns and a symmetric positive definite matrix is a polar decomposition, so $\polar(G)=O_0$ and the
next iterate is on the same line. A combination with non-negative weights of gradients $O_0(Y_0-\tau_s\Sigma)$,
$\tau_s<\tau_\ast$, has the same polar factor. For the flow, $\partial\psi(G)=\{O_0\}$ at these full-rank gradients, so
the segment solves the inclusion, and the solution is unique.
(b) $\tau_\ast=\min_x\frac{x^\top Y_0x}{x^\top\Sigma x}\le\frac{\tr Y_0}{\tr\Sigma}=\tau_{\mathrm{ls}}$, with
equality iff $\Sigma^{-1/2}Y_0\Sigma^{-1/2}$ is a multiple of the identity, i.e.\ $\Sigma E_0^\top E_0\Sigma=c^2\Sigma^2$;
and $\tau_{\mathrm{ls}}=\nnuc{E_0\Sigma}/\tr\Sigma\le\nop{E_0}$. The eigenvalues of $\Sigma^{-1/2}Y_0\Sigma^{-1/2}$ are
positive with product $\det Y_0/\det\Sigma=\det\abs{E_0}$, so the smallest is at most the geometric mean
$(\prod_is_i(E_0))^{1/d}$. Finally $\Sigma\abs{E_0}^2\Sigma\succeq s_d(E_0)^2\Sigma^2$ and the square root is operator
monotone, so $Y_0\succeq s_d(E_0)\Sigma$.
(c) $\polar(O_0M)=O_0\polar(M)$, and $\tr M_1=\tr Y_0-\tau_{\mathrm{ls}}\tr\Sigma=0$. Write $M_1=M_+-M_-$ (positive and
negative parts, with orthogonal ranges) and $S=\tr M_+=\tr M_-$, so that $\nnuc{G_1}=\tr\abs{M_1}=2S$ and
$2(F(W_1)-F_\ast)=\tr(\Sigma^{-1}M_1^2)=\tr(\Sigma^{-1}M_+^2)+\tr(\Sigma^{-1}M_-^2)$. The function
$M\mapsto\tr(\Sigma^{-1}M^2)$ is convex, so on $\{M\succeq0,\ \tr M=S,\ \mathrm{range}(M)\subset\mathrm{range}(M_+)\}$ it is
maximal at an extreme point $Sxx^\top$: $\tr(\Sigma^{-1}M_+^2)\le S^2x^\top\Sigma^{-1}x$, and similarly
$\tr(\Sigma^{-1}M_-^2)\le S^2y^\top\Sigma^{-1}y$ with unit vectors $x\perp y$. By Ky Fan's inequality
$x^\top\Sigma^{-1}x+y^\top\Sigma^{-1}y\le1/\lambda_d+1/\lambda_{d-1}$. Since $J\le\tr\Sigma$ at every iterate,
$F(W_2)-F_\ast\le F(W_1)-F_\ast-\nnuc{G_1}^2/(2\tr\Sigma)$, and
$\nnuc{G_1}^2/(2\tr\Sigma(F(W_1)-F_\ast))\ge4S^2/\big(\tr\Sigma\,S^2(1/\lambda_d+1/\lambda_{d-1})\big)
=4p_dp_{d-1}/(p_d+p_{d-1})\ge2p_d$.
\end{proof}

The factorization, the value of $F$ along the line, $\tau_{\mathrm{ls}}$, the rank-drop times and the trace-zero
identity are (K2). Added here: the frozen polar factor for every schedule, for momentum and for the flow, the
inequality $\tau_{\mathrm{ls}}\ge\tau_\ast$ and the bounds on $\tau_\ast$ (the bounds are due to the mathematical referee
of this section), and the second-step bound (c), also due to that referee: it extends
Theorem~\ref{thm:mad}(b) at $t=1$ to non-commuting starts (for $d=2$ it is the exact rate of (K2), valid at every
step). After the second step $M_t$ is no longer symmetric in general and no such bound is known. The bounds on
$\tau_\ast$ and the second-step bound were added after the literature checks and have not been searched for.

\begin{proposition}[rotation between the error and its gradient]\label{prop:Q} \textbf{[P]}
Let $E$ have rank $d$. Then
\[
\polar(E\Sigma)=\polar(E)\,Q,\qquad Q=\polar(\abs E\Sigma)=\abs E\Sigma\,(\Sigma\abs E^2\Sigma)^{-1/2}\in SO(d),
\]
and $Q=I$ if and only if $\abs E$ commutes with $\Sigma$. Along the flow $\dot E=-\polar(E\Sigma)$, as long as $E$ has
rank $d$ (and, for (c), simple singular values $s_i$ with right singular vectors $w_i$):
\begin{enumerate}[label=(\alph*)]
\item $\frac d{dt}F=-\nnuc G$ and $\frac d{dt}\nnuc G=-\tr\Sigma$ (Proposition~\ref{prop:flow} and \eqref{eq:Jconst});
\item $\frac d{dt}\frac12\nF E^2=-\tr(\Sigma^{-1}\abs G)<0$ and $\frac d{dt}\nnuc E=-\tr Q$;
\item $\frac d{dt}s_i=-w_i^\top Qw_i$. If $d=2$, $Q$ is the rotation of angle $\varphi$ with
$\cos\varphi=\tr(\abs E\Sigma)/\nnuc{\abs E\Sigma}\in(0,1]$ and both singular values of the error decrease at speed
$\cos\varphi$; explicitly $\varphi=\arg\big(1+\rho_E\rho_\Sigma e^{2i\theta}\big)$ with
$\rho_E=\frac{s_1-s_2}{s_1+s_2}$, $\rho_\Sigma=\frac{\lambda_1-\lambda_2}{\lambda_1+\lambda_2}$ and $\theta$ the angle
between the eigenbases of $\abs E$ and $\Sigma$, so that $\cos\varphi\ge\sqrt{1-\rho_E^2\rho_\Sigma^2}$;
\item (conservation law, discrete time) if $K=d=2$, every polar step $E^+=E-\eta\,\polar(E\Sigma)$, of any length,
conserves $s_1(E)-s_2(E)$ as long as $\det E$ keeps its sign;
\item if $d\ge3$, $w^\top Qw$ can be negative for a top singular vector $w$ of $E$: \emph{the spectral norm of the error
can increase along the flow}, while $F$ and $\nF E$ decrease. Example with integer entries
(\texttt{counterexample\_opnorm\_integer.py}): $E_0=\diag(300,200,1)$,
$\Sigma=\left(\begin{smallmatrix}2&4&12\\4&9&20\\12&20&89\end{smallmatrix}\right)$ ($\det\Sigma=2$): $w=e_1$ and
$w^\top Qw=Q_{11}=-0.01433$; $\nop{E(t)}$ increases from $300$ to $300.86$ during the first phase
($t\le\tau_\ast=14.34$) and reaches $303.27$ at $t=33.5$, while $\nF{E(t)}$ decreases, and the flow ends at
$T=1.20\,\nop{E_0}$;
\item (the increase needs an ill-conditioned error) if $d=3$, with $g=(s_1s_2s_3)^{1/3}$ and $\{i,j,k\}=\{1,2,3\}$,
\[
w_i^\top Qw_i\ \ge\ \frac{3g(s_i+g)-s_i(s_j+s_k)}{(s_i+s_j)(s_i+s_k)} ;
\]
in particular $\frac d{dt}\nop E<0$ for every $\Sigma$ as soon as $s_1^2<(2+\sqrt3)^3s_2s_3\approx52\,s_2s_3$, hence as
soon as the condition number of $E$ is below $(2+\sqrt3)^{3/2}\approx7.2$.
\end{enumerate}
\end{proposition}
\begin{proof}
$Q^\top Q=(\Sigma\abs E^2\Sigma)^{-1/2}\Sigma\abs E^2\Sigma(\Sigma\abs E^2\Sigma)^{-1/2}=I$, and
$\polar(E)Q=E\abs E^{-1}\abs E\Sigma\abs G^{-1}=G\abs G^{-1}$ with $\abs G=(\Sigma\abs E^2\Sigma)^{1/2}$;
$\det Q>0$ since $\det(\abs E\Sigma)>0$. $Q=I$ iff $\abs E\Sigma$ is symmetric positive definite iff $\abs E\Sigma=\Sigma\abs E$.
(b) $\frac d{dt}\frac12\nF E^2=-\inner E{\polar(E)Q}=-\tr(\abs EQ)$
and $\abs EQ=\abs E^2\Sigma\abs G^{-1}=\Sigma^{-1}\abs G^2\Sigma^{-1}\Sigma\abs G^{-1}=\Sigma^{-1}\abs G$, whose trace is
$\tr(\Sigma^{-1/2}\abs G\Sigma^{-1/2})>0$; $\frac d{dt}\nnuc E=-\inner{\polar(E)}{\polar(E)Q}=-\tr Q$.
(c) $\dot s_i=u_i^\top\dot Ew_i=-u_i^\top\polar(E)Qw_i=-w_i^\top Qw_i$. For $d=2$, $Q^\top\abs E\Sigma=\abs G$ is symmetric,
and for a rotation $R_\varphi$ and a symmetric $Y$, $\tr(R_\varphi Y)=\cos\varphi\,\tr Y$; hence
$\tr(\abs E\Sigma)=\cos\varphi\,\tr\abs G$, and $\tr(\abs E\Sigma)=\tr(\abs E^{1/2}\Sigma\abs E^{1/2})>0$. Writing
$\abs E=\frac{s_1+s_2}2(I+\rho_ES_a)$ and $\Sigma=\frac{\lambda_1+\lambda_2}2(I+\rho_\Sigma S_b)$ with symmetric
reflections $S_a,S_b$ of axes at angles $a,b$, the product $\abs E\Sigma$ is proportional to
$I+\rho_E\rho_\Sigma S_aS_b+\rho_ES_a+\rho_\Sigma S_b$, where $S_aS_b$ is the rotation of angle $2(a-b)$; the polar factor
of a $2\times2$ matrix of positive determinant is its normalized component in the plane of rotations (see (d)), here
the rotation of angle $\arg(1+\rho_E\rho_\Sigma e^{2i(a-b)})$, and $\abs{\arg(1+xe^{i\alpha})}\le\arcsin x$ for $0\le x<1$.
(d) Write a $2\times2$ matrix as $M=M_c+M_a$ with $M_c$ in the span of the rotations (multiples of rotations) and $M_a$
in the span of the reflections; these planes are orthogonal, and with $m_c=\nF{M_c}/\sqrt2$, $m_a=\nF{M_a}/\sqrt2$ the
singular values are $m_c+m_a$ and $\abs{m_c-m_a}$, and $\det M=m_c^2-m_a^2$. So $s_1-s_2=2\min(m_c,m_a)$. If
$\det E>0$ then $\det(E\Sigma)>0$ and $\polar(E\Sigma)$ is a rotation: the step changes only $E_c$, and
$s_1-s_2=2m_a$ before and after if the determinant is still positive. If $\det E<0$ the polar factor is a reflection
and $m_c$ is conserved.
(f) Let $B=Q^\top\abs E=\abs G\Sigma^{-1}$. It is similar to $\Sigma^{-1/2}\abs G\Sigma^{-1/2}$, so its eigenvalues
$\mu_1,\mu_2,\mu_3$ are positive, with product $\det\abs E=g^3$; by the inequality between arithmetic and geometric
means, $\tr B=\sum_i\mu_i\ge3g$ and the sum $e_2(B)$ of its principal $2\times2$ minors is $\sum_{i<j}\mu_i\mu_j\ge3g^2$. In
the basis $(w_i)$, $B_{ij}=s_jQ_{ji}$ with $Q_{ji}=w_j^\top Qw_i$, so $\tr B=\sum_is_iq_i$ with $q_i=Q_{ii}$, and, a
rotation of $\R^3$ being equal to its cofactor matrix, $e_2(B)=s_2s_3q_1+s_1s_3q_2+s_1s_2q_3$. A rotation of angle $t$
about the unit axis $n$ has $q_i=\cos t+(1-\cos t)n_i^2$, hence $q_j+q_k\le1+q_i$. Therefore
\[
3g(s_i+g)\le s_i\tr B+e_2(B)=q_i(s_i^2+s_js_k)+(q_j+q_k)s_i(s_j+s_k)\le q_i(s_i+s_j)(s_i+s_k)+s_i(s_j+s_k).
\]
For $i=1$, with $a=s_2/s_1$, $b=s_3/s_1$ and $h=(ab)^{1/3}$: $a+b\le1+ab=1+h^3$ and
$3h+3h^2-1-h^3=-(h+1)(h^2-4h+1)>0$ for $2-\sqrt3<h\le1$, i.e.\ for $s_1^2<(2+\sqrt3)^3s_2s_3$; and
$s_1^2/(s_2s_3)\le(s_1/s_3)^2$.
\end{proof}

Item (d) was found by the literature referee of this section while checking (c); it explains the observation of Shulgin
et al.\ (Figure~13) that in the $2\times2$ case the singular values move along a $45^\circ$ line. Item (e) has an
$\ell_\infty$ analogue: for a ball-proximal method in the $\ell_\infty$ norm on a two-dimensional quadratic the distance
to the solution in that norm can increase (Gruntkowska and Richt\'arik, arXiv:2510.00823, Figure~1). The ``if'' part of
$Q=I$ is in Shulgin et al.\ (Corollary~4), and as the ``shared frames'' condition of Beneventano, Abdelmoneum and
Poggio (arXiv:2606.08388, Definition~1), whose Theorem~1(i) gives a formula of the same type as (c) for singular values
along a projected polar flow in underdetermined regression. The factor $Q$, items (b), (d), (e) and the formula
$\dot s_i=-w_i^\top Qw_i$ of (c) were not found; the angle formula of (c) and item (f) were added during the revision by
the mathematical referee of this section and have not been searched for. The example of (e) is also due to that
referee, who certified the sign of $Q_{11}$ in exact rational
arithmetic; the example replaces a numerical one of the first version (condition numbers $10^3$,
$w^\top Qw=-0.0958$, \texttt{counterexample\_opnorm.py}) in which the top singular value of the error was double, so
that (c) did not apply. By (f) no well-conditioned example exists in dimension three; in the referee's searches
\textbf{[N]} a negative value also required a condition number of $\Sigma$ of several hundreds (the smallest value of
$w_1^\top Qw_1$ found with both condition numbers below $c$ was $+0.39$, $+0.02$, $-0.09$ for $c=50,400,1000$).

\begin{theorem}[nuclear-scaled step in the commuting case: contraction after the first step]\label{thm:mad}
\textbf{[P]; the map is (K3)}
In the commuting case let $a(t)\in(0,\infty)^d$ be the amplitudes along the nuclear-scaled step with $L=\tr\Sigma$, and
let $\E_p$, $\Var_p$, $\MAD_p$ denote the mean, the variance and the mean absolute deviation around the mean under the
weights $p$. As long as all amplitudes are positive, $a(t+1)=\abs{a(t)-\E_p[a(t)]}$ and:
\begin{enumerate}[label=(\alph*)]
\item $\displaystyle\frac{F(W_{t+1})-F_\ast}{F(W_t)-F_\ast}=\frac{\Var_p(a(t))}{\E_p[a(t)^2]}$, which equals
$1-\MAD_p(a(t-1))^2/\Var_p(a(t-1))$ for $t\ge1$;
\item for $t\ge1$ this factor is at most $1-\frac{4p_{(1)}p_{(2)}}{p_{(1)}+p_{(2)}}\le1-\frac2{\kop}$, where
$p_{(1)}\le p_{(2)}$ are the two smallest weights: after the first step the exponent is at least twice the certified
one. The constant is sharp: it is a supremum over positive amplitudes, attained iff $d=2$ (where the factor is
$(p_1-p_2)^2$ for every $t\ge1$). At $t=0$ the supremum of the factor is the certified $1-1/\kop$ (error almost
concentrated on the flattest direction, the other amplitudes small and positive); it is not attained, and for an error
exactly of rank one on $v_d$ the factor is $(1-1/\kop)^2$;
\item conversely, every invariant ray of the map (amplitudes $a\in(0,\infty)^d$ with $a(t+1)=\mu a(t)$, $\mu>0$) is
two-valued, $a=\delta((1-\pi)\ones_A+\pi\ones_{A^c})$ with $\pi=p(A)<\frac12$ and $\mu=1-2\pi$; the only other
positive configuration mapped to a multiple of itself is the constant one ($\mu=0$: convergence in one step). On such a
ray with $A=\{k\}$ the gradient has nuclear rank in $[2,4(1-1/d)]$, equal to $2$ iff $d=2$. For $d\ge3$ these rays are
linearly unstable: the linearized map has eigenvalues $1-2\pi$, $0$, $+1$ ($\abs A-1$ times) and $-1$
($\abs{A^c}-1$ times), so perturbations grow like $(1-2\pi)^{-t}$ relative to the ray.
\end{enumerate}
\end{theorem}
\begin{proof}
Every amplitude moves by $\eta=\nnuc G/\tr\Sigma=\E_p[a]$, and $F-F_\ast=\frac{\tr\Sigma}2\E_p[a^2]$. (a)
$\E_p[a(t+1)^2]=\Var_p(a(t))$ and $\E_p[a(t)]=\E_p\abs{a(t-1)-\E_pa(t-1)}$. (b) Let $X$ be centred under $p$ and
$S=\E_p[X_+]=\E_p[X_-]$, so $\E_p\abs X=2S$. On $\{X>0\}$, $\sum p_kx_k^2$ is a convex function of $x\ge0$ on
$\{\sum p_kx_k=S\}$, maximal at a vertex: it is at most $S^2/\min_{X_k>0}p_k$; similarly on $\{X<0\}$. The two minima are
attained at distinct indices, so $\E_pX^2\le S^2(1/p_{(1)}+1/p_{(2)})$ and
$(\E_p\abs X)^2/\E_pX^2\ge4p_{(1)}p_{(2)}/(p_{(1)}+p_{(2)})\ge2p_{(1)}$, with equality only for
$X=S(e_{(1)}/p_{(1)}-e_{(2)}/p_{(2)})$, which corresponds to amplitudes with zero entries when $d>2$ (there the true
iteration keeps the zeros and contracts faster). At $t=0$, $\Var_p(a)/\E_p[a^2]=1-(\E_pa)^2/\E_p[a^2]\to1-p_d$ when
$a\to e_d$, while at $a=e_d$ the step multiplies the only amplitude by $1-p_d$. (c) If $a(t+1)=\mu a(t)$ with signs
$s_k=\sign(a_k-m)$, $m=\E_pa>0$, then $a_k(1-\mu s_k)=m$ when $\mu>0$ (no $s_k$ vanishes, and $\mu\ne1$), so $a$ takes two
values; with $\pi$ the weight of $\{s_k=+1\}$, the constraint $m=\E_pa$ gives $\mu=1-2\pi$, and consistency of the signs
requires $\pi<\frac12$; $\mu=0$ means $a_k=m$ for all $k$. On the ray with $A=\{k\}$,
$\nnuc G=2\pi(1-\pi)\delta\tr\Sigma$ and
$\nF G^2=\delta^2\big((1-\pi)^2\lambda_k^2+\pi^2\sum_{j\ne k}\lambda_j^2\big)$ with $\lambda_k=\pi\tr\Sigma$, so
$\nr(G)=4(1-\pi)^2/\big((1-\pi)^2+\sum_{j\ne k}p_j^2\big)$ with
$(1-\pi)^2/(d-1)\le\sum_{j\ne k}p_j^2\le(1-\pi)^2$, the upper bound being attained iff $d=2$. The map is differentiable
at a point of the ray, with derivative $\diag(s)(I-\ones p^\top)$: it is $+1$ on vectors supported in $A$ with zero
$p$-mean, $-1$ on those supported in $A^c$, and on the span of $\ones_A,\ones_{A^c}$ it has trace $1-2\pi$ and
determinant $0$.
\end{proof}

For $d=2$ the bound of (b) is the exact rate of (K2), and the statement at $t=0$ is the tightness statement of Shulgin
et al.\ (Theorem~3 and Appendix~E.2). By Proposition~\ref{prop:phase1}(c) the bound of (b) at $t=1$ holds for every
full-rank start, commuting or not. The form $1-\MAD^2/\Var$, the bound in general dimension and the classification of
(c) were not found; the instability, the range of the nuclear rank and the corrections of (b) are due to the
mathematical referee of this section, were added after the literature checks and have not been searched for.

\begin{proposition}[extinction time of the flow]\label{prop:ext} \textbf{[P]; (a) and the bounds of (b) are [K]}
For the flow $\dot E\in-\partial\psi(E\Sigma)\,$:
\begin{enumerate}[label=(\alph*)]
\item in the commuting case $T=\nop{E_0}$ whatever $\Sigma$ (soft-thresholding path of (K3)): the lower bound of
\eqref{eq:sandwich} is attained;
\item if $E_0=ue_0^\top$ has rank one ($\norm u=1$), $E(t)=ue(t)^\top$ with $\dot e=-\Sigma e/\norm{\Sigma e}$, the
gradient flow of $\frac12e^\top\Sigma e$ parametrized by arc length, and $T$ is the length of that curve:
$\norm{e_0}_2\le T\le\min\{\sum_k\abs{v_k^\top e_0},\ (1+2.5\sqrt{\log\kappa(\Sigma)})\norm{e_0}_2\}$, and there are
instances with $T\ge\min\{0.7\sqrt d,0.45\sqrt{\log\kappa(\Sigma)}\}\norm{e_0}_2$;
\item in general $\nop{E_0}\le T\le\sqrt{2(F(W_0)-F_\ast)/\lambda_d}\le\sqrt{\tr\Sigma/\lambda_d}\,\nop{E_0}$;
\item if $d=2$ and $E_0$ has singular values $s_1\ge s_2$, then
$T\le\sqrt{s_1s_2}+\sqrt2\,(s_1-s_2)\le\frac{9\sqrt2}8\nop{E_0}\approx1.59\,\nop{E_0}$, whatever $\Sigma$.
\end{enumerate}
\end{proposition}
\begin{proof}
(a) The curve $\sum_k(a_k-t)_+u_kv_k^\top$ is absolutely continuous with derivative
$-\sum_{a_k>t}u_kv_k^\top=-\polar(G(t))$ for almost every $t$, so it is the solution, by uniqueness
(Proposition~\ref{prop:flow}). (b) $G=u(\Sigma e)^\top$, $\polar(G)=u(\Sigma e)^\top/\norm{\Sigma e}$, so this curve
satisfies the inclusion and is the solution, by uniqueness; up to its parametrization it is
$s\mapsto e^{-s\Sigma}e_0$, each coordinate $v_k^\top e$ is monotone along it, so the length is at most
$\sum_k\abs{v_k^\top e_0}$; the other bounds are those of Gupta, Balakrishnan and Ramdas for path lengths of gradient
flows of quadratics (JMLR 2021, Theorems 10 and 18). (c) is \eqref{eq:sandwich}. (d) If $s_1=s_2$ the line search is
exact and $T=s_1$. Otherwise, during the full-rank phase both singular values decrease at the same speed
(Proposition~\ref{prop:Q}(c)), so $s_1-s_2$ is conserved and the phase ends when the smaller one vanishes, at
$\tau_\ast\le\sqrt{s_1s_2}$ (Proposition~\ref{prop:phase1}(b)); the error then has rank one and norm $s_1-s_2$, and by
(b) the remaining time is at most $\sqrt2(s_1-s_2)$. The maximum of $\sqrt x+\sqrt2(1-x)$ on $[0,1]$ is at $x=1/8$.
\end{proof}

Item (d), due to the mathematical referee of this section, answers in dimension two a question left open in the first
version (is $T\le C\nop{E_0}$ with $C$ independent of $\Sigma$?); it was added after the literature checks and has not
been searched for. By (b) such a constant must grow at least like
$\min\{\sqrt d,\sqrt{\log\kappa(\Sigma)}\}$; for $d=2$ the computed supremum is $\sqrt2$, approached in the rank-one
limit (largest value $1.4138$ over $3000$ instances with $\kappa(\Sigma)\le10^6$ in \texttt{check\_ls\_additions.py}),
and for rank-one errors with widely spread spectra the referees found $T/\nop{E_0}=1.72$, $2.43$, $4.00$ for
$d=3,6,16$, which suggests $C=\sqrt d$ \textbf{[N]}.

\noindent\emph{Numerical support} (scripts of Appendix~\ref{app:checks}): in the commuting case the simulated flow follows the soft-thresholding path within
the step for $\kappa(\Sigma)$ from $1$ to $10^4$; on generic instances ($6\times6$ and $10\times4$,
$\kappa(\Sigma)\le400$) $T/\nop{E_0}$ is between $1.00$ and $1.24$ while $\sqrt{2D_0/\muop}/\nop{E_0}$ is between $1.6$
and $16.6$ and $\sqrt{\kop}$ up to $23.9$ (the referees' runs: $0.99$ to $1.36$ over $48$ instances, and $1.00$ to
$1.35$ for Gaussian starts with $d\in\{2,3,6\}$ and $\kappa(\Sigma)\le64$); for $d=5$ and
$\kappa(\Sigma)=1,10^2,10^4$ from the same generic start, $T/\nop{E_0}=1.00,1.14,1.23$ (Figure~\ref{fig:ls}(b)).

\begin{figure}[t]
\centering
\includegraphics[width=\textwidth]{figT1_ls_model}
\caption{The least-squares model. (a) Polar flow from an error that commutes with $\Sigma$: the amplitudes follow the
soft-thresholding path $(a_k-t)_+$ for both condition numbers (published; (K3)). (b) Generic start: the spectral norm
of the error for three condition numbers; the extinction time stays within $25\%$ of $\nop{E_0}$. (c) Discrete steps on
$K=d=100$, spectrum $1/k$ ($\kF=100$, $\kop=519$): the nuclear-scaled polar step is ahead of gradient descent for the
first hundreds of iterations and behind afterwards. Script: \texttt{make\_paper\_figures.py}.}\label{fig:ls}
\end{figure}

\begin{remark}[head and tail; link with B2]\label{rem:headtail}
The model separates two regimes, both described in the literature ((K3), (K4)). \emph{Head:} the polar step shrinks
all singular values of the error at the same speed, independently of the curvature (exactly in the commuting case).
\emph{Tail:} once small singular values of the error have come close to zero, every further polar step moves them by
the full length of the step (the absolute value in the map $a\mapsto\abs{a-\E_pa}$). Three exact reference rates per
step for $F-F_\ast$ are available: $(1-2p_d)^2$ on the invariant ray of the flattest direction (exponent
$\approx4\lambda_d/\tr\Sigma$), $(1-p_d)^2$ for an error exactly of rank one on $v_d$ (exponent
$\approx2\lambda_d/\tr\Sigma$; on a rank-one error $ue^\top$ the nuclear-scaled step is gradient descent with step
$1/\tr\Sigma$ on the row, $e^+=e-\Sigma e/\tr\Sigma$), and the certified $1-p_d$. Gradient descent with step $1/\lambda_1$ has exponent $2\lambda_d/\lambda_1$ in its tail: \emph{compared with
the flattest ray, the nuclear-scaled polar step is slower than gradient descent by the factor
$\tr\Sigma/(2\lambda_1)$}, half the stable rank of the data matrix (Proposition~\ref{prop:struct}(b)), and faster when
that stable rank is below $2$. The factor $(1-2w)^2$ is Corollary~5 of Shulgin et al., who compare the resulting rate
with gradient descent for a spectrum with one small eigenvalue (text after their Theorem~3 and Appendix~E.2; comment
(i) after the corollary).

Which of these rates a run follows is a numerical matter \textbf{[N]}, because the rays are unstable
(Theorem~\ref{thm:mad}(c)). From generic, non-commuting starts the tail is close to the flattest ray: in the
referee's runs ($K=d=20$, five seeds) the exponent per step is $4.12\,\lambda_d/\tr\Sigma$ for the spectrum $1/k$ and
$4.01$ to $4.02$ for $1/k^2$, the top right singular vector of the error is aligned with $v_d$ (cosines $0.996$ to
$1.000$) and the nuclear rank of the gradient is $3.3$ and $2.5$ on average; the ratio of the exponents of gradient descent and of
the polar step is $1.79$ and $0.80$, against $\tr\Sigma/(2\lambda_1)=1.80$ and $0.80$: the polar tail is slower for the
first spectrum and faster for the second. In the commuting case there is no such alignment in exact arithmetic: within
the first few iterations the dominant amplitude settles on an index $k^\ast$ determined by the initial amplitudes and
the weights (over $40$ random starts with $d=20$, it is the index of the largest initial amplitude in $25$ runs for the
spectrum $1/k$ and $23$ runs for $1/k^2$, and it no longer changes after the eighth step;
\texttt{check\_ls\_additions.py}, block (9), $50$-digit arithmetic), and the exponent of $F$ is then close to
$-2\log(1-2p_{k^\ast})\approx4\lambda_{k^\ast}/\tr\Sigma$ (between $1.00$ and $1.26$ times this value over steps $200$ to
$400$), so the flattest ray is the worst case; in double precision the matrix iteration follows
the scalar map for thousands of iterations, then leaves the commuting manifold and joins the generic regime
(similarly, at a rank-deficient commuting iterate a perturbation $\epsilon u_iv_j^\top$, with $\sigma_i\ne0$ and
$\sigma_j=0$, is multiplied to first order by $1-\eta\lambda_j/(\abs{\sigma_i}\lambda_i)$ at each step, of modulus larger
than one as soon as $\eta\lambda_j>2\abs{\sigma_i}\lambda_i$: exact zeros are not robust in floating-point runs). For
the run of Figure~\ref{fig:ls}(c) only the exponent can be compared: for $K=d=100$ and spectrum $1/k$
($\tr\Sigma/\lambda_1=5.19$), the relative objective gap after $10$, $100$ and $1000$ iterations is
$4.0\cdot10^{-3}$, $1.3\cdot10^{-4}$, $4.5\cdot10^{-8}$ for the nuclear-scaled polar step, $0.23$, $9.0\cdot10^{-3}$,
$1.8\cdot10^{-11}$ for gradient descent with step $1/\lambda_1$ and $0.14$, $6.0\cdot10^{-4}$, $3.3\cdot10^{-21}$ with
exact line search; over the last $500$ of $3000$ iterations the exponent per step of the polar run is $0.00809$
($0.00833$ and $0.00817$ on the two preceding windows), close to that of the flattest invariant ray, $0.00773$ (the
bound of Theorem~\ref{thm:mad}(b) is $0.00386$), against $0.0204$ for gradient descent ($2.5$ times larger, for
$\tr\Sigma/(2\lambda_1)=2.6$). But the rays of the ten flattest directions have exponents between $0.0077$ and
$0.0085$, and at iteration $3000$ the error, although essentially of rank one ($s_1^2$ carries $99.9\%$ of
$\nF E^2$), is not aligned with $v_d$ (cosine $0.54$; nuclear rank of the gradient $3.6$): the bottom of this spectrum
is nearly degenerate, and this run does not discriminate the flattest ray from its neighbours. The alignment of the
error with the flattest eigenvector in the tail, where it is observed (the $K=d=20$ runs above), is the counterpart of
the alignment proved for gradient descent in the note B2 (Proposition B2-gd); here it is not proved.
\end{remark}

\section{Stochastic polar steps: smoothing by the noise and the nuclear median}\label{sec:noise}

Let $\tilde g=G+\xi$ be a stochastic gradient, $G$ deterministic and $\xi$ a random matrix. Define the
\emph{noise-smoothed nuclear norm} and the mean polar field
\[
\psi_\xi(G)=\E\big[\nnuc{G+\xi}-\nnuc\xi\big],\qquad \Phi_\xi(G)=\E\,\polar(G+\xi).
\]
No moment is needed: the integrand of $\psi_\xi$ is bounded by $\nnuc G$, and $\polar$ is a bounded Borel map (the
pointwise limit of $M(M^\top M+\epsilon I)^{-1/2}$). When $\E\nnuc\xi<\infty$, $\psi_\xi(G)=\E\nnuc{G+\xi}-\E\nnuc\xi$.

\begin{theorem}[the expected polar step is a subgradient]\label{thm:smooth} \textbf{[K]}
$\psi_\xi$ is convex, $\abs{\psi_\xi(G)-\psi_\xi(G')}\le\nnuc{G-G'}$, and
\[
\Phi_\xi(G)\in\partial\psi_\xi(G);\qquad \Phi_\xi(G)=\nabla\psi_\xi(G)\ \text{ if }\ \Prob(\operatorname{rank}(G+\xi)<r)=0 .
\]
In particular $\Phi_\xi$ is monotone, and $\inner{G-G'}{\Phi_\xi(G)}\ge\psi_\xi(G)-\psi_\xi(G')$. If $\E\nnuc\xi<\infty$, the
same holds for the nuclear-scaled map with $h_\xi(G)=\frac12\E[\nnuc{G+\xi}^2-\nnuc\xi^2]$:
$\E\,\sharpop(G+\xi)\in\partial h_\xi(G)$.
\end{theorem}
\begin{proof}
For each realization $\polar(G+\xi)\in\partial\psi(G+\xi)$:
$\nnuc{G'+\xi}\ge\nnuc{G+\xi}+\inner{\polar(G+\xi)}{G'-G}$; take expectations. If $G+\xi$ has full rank almost surely,
$\psi$ is differentiable at $G+\xi$ almost surely and dominated convergence (difference quotients are bounded by
$\nnuc{\Delta}$) shows that the directional derivatives of $\psi_\xi$ at $G$ are linear, so the convex function
$\psi_\xi$ is differentiable at $G$. For $h_\xi$ use
$\frac12\psi(M')^2\ge\frac12\psi(M)^2+\psi(M)\inner O{M'-M}$ for $O\in\partial\psi(M)$.
\end{proof}

This is the interchange of expectation and subdifferential for a convex integrand, i.e.\ the gradient formula of
randomized smoothing (Duchi, Bartlett and Wainwright, arXiv:1103.4296; Abernethy et al., COLT 2014, quoting Bertsekas
1973; for the support function of a set, Dubois-Taine et al., NeurIPS 2022, Proposition~2), here for the spectral unit
ball. For Gaussian matrix perturbations injected by the algorithm it is written verbatim by Jiang, Mhammedi, Mohri and
Mokhtari (arXiv:2602.08232, Eq.~(6)--(7)); that gradient noise smooths the polar map is stated by Xie (arXiv:2609.32861,
Proposition~3.1), Jacobs and Burkholz (arXiv:2610.16868) and Chen, Li, Wang and Lu (arXiv:2610.31728). What is added below
is the combination with Proposition~\ref{prop:dual} for the method's own noise.

\begin{theorem}[stochastic polar steps are stochastic mirror descent]\label{thm:smd} \textbf{[P]}
Assume \eqref{eq:consts} with $\muop>0$, let $\xi_t$ be i.i.d.\ copies of $\xi$, independent of the past (additive
noise), and $W_{t+1}=W_t-\eta_t\,\polar(\nabla F(W_t)+\xi_t)$, $G_t=\nabla F(W_t)$. The set $\arg\min\psi_\xi$ is nonempty,
convex and compact. For $G^\dagger\in\arg\min\psi_\xi$ and $W^\dagger=\nabla F^\ast(G^\dagger)$:
\begin{enumerate}[label=(\alph*)]
\item $\displaystyle\E\big[D_F(W_{t+1},W^\dagger)\,\big|\,W_t\big]\le D_F(W_t,W^\dagger)-\eta_t\big(\psi_\xi(G_t)-\min\psi_\xi\big)+\frac{\Lop}2\eta_t^2$;
\item for any lengths, with $\bar G_T=\sum_{t<T}\eta_tG_t/\sum_{t<T}\eta_t$,
\[
\E\,\psi_\xi(\bar G_T)-\min\psi_\xi\ \le\ \frac{D_F(W_0,W^\dagger)+\frac{\Lop}2\sum_{t<T}\eta_t^2}{\sum_{t<T}\eta_t};
\]
\item if $\sum_t\eta_t=\infty$ and $\sum_t\eta_t^2<\infty$, then almost surely $G_t$ converges to a point of
$\arg\min\psi_\xi$;
\item consequently, for such lengths, $W_t\to W_\ast$ almost surely if $\arg\min\psi_\xi=\{0\}$, and if
$0\notin\arg\min\psi_\xi$ then almost surely
$\lim_t F(W_t)-F_\ast\ge\frac1{2\Lop}\min_{G\in\arg\min\psi_\xi}\nnuc G^2>0$, \emph{whatever the schedule in this
class}.
\end{enumerate}
No variance assumption is needed: the steps are bounded.
\end{theorem}
\begin{proof}
$\psi_\xi$ is continuous and coercive: with $R$ such that $\Prob(\nnuc\xi\le R)>\frac12$,
$\psi_\xi(G)\ge\nnuc G(2\Prob(\nnuc\xi\le R)-1)-2R$. (a) With $O_t=\polar(G_t+\xi_t)$, the three-point identity of mirror
descent with mirror map $F^\ast$ reads
$D_{F^\ast}(G^\dagger,G_{t+1})=D_{F^\ast}(G^\dagger,G_t)-\eta_t\inner{O_t}{G_t-G^\dagger}+D_{F^\ast}(G_t,G_{t+1})$. By
Proposition~\ref{prop:dual}(b), $D_{F^\ast}(G^\dagger,G_t)=D_F(W_t,W^\dagger)$ and
$D_{F^\ast}(G_t,G_{t+1})=D_F(W_{t+1},W_t)\le\frac{\Lop}2\nop{W_{t+1}-W_t}^2\le\frac\Lop2\eta_t^2$. Conditionally on $W_t$,
$\E\inner{O_t}{G_t-G^\dagger}=\inner{\Phi_\xi(G_t)}{G_t-G^\dagger}\ge\psi_\xi(G_t)-\psi_\xi(G^\dagger)$ by
Theorem~\ref{thm:smooth}. (b) Sum, and use the convexity of $\psi_\xi$. (c) By the Robbins--Siegmund lemma,
$D_F(W_t,W^\dagger)$ converges and $\sum_t\eta_t(\psi_\xi(G_t)-\min\psi_\xi)<\infty$ almost surely; $W_t$ is bounded
($D_F(W_t,W^\dagger)\ge\frac{\muop}2\nop{W_t-W^\dagger}^2$), a cluster point of $G_t$ is a minimizer, and the convergence of
$D_F(W_t,W')$ for that minimizer (simultaneously for a countable dense set of minimizers) gives the convergence of the
whole sequence. (d) $F(W_t)-F_\ast\ge\nnuc{G_t}^2/(2\Lop)$ by Proposition~\ref{prop:dual}(c).
\end{proof}

\begin{theorem}[nuclear medians and zeros of the mean field]\label{thm:median} \textbf{[P]; scalar and sign versions are [K]}
\begin{enumerate}[label=(\alph*)]
\item $\arg\min\psi_\xi=-\Med(\xi)$, where $\Med(\xi)=\arg\min_C\E[\nnuc{\xi-C}-\nnuc\xi]$ is the set of \emph{nuclear
medians} of $\xi$ (the median of a law in a normed space, for the nuclear norm). Every zero of the mean field is a
nuclear median: $\Phi_\xi(G)=0\Rightarrow G\in-\Med(\xi)$. The converse holds if $\Prob(\operatorname{rank}(G+\xi)<r)=0$,
and can fail otherwise: for $\xi=\pm\frac12\diag(1,-1)$ with probability $\frac12$, $G=\frac12I$ minimizes $\psi_\xi$ and
$\Phi_\xi(G)=\frac12I$; for the scalar law $\Prob(\xi=1)=0.4$, $\Prob(\xi=-1)=0.6$, the minimizer is $g=1$ and
$\abs{\Phi_\xi}\ge0.2$ everywhere.
\item If $\xi$ and $-\xi$ have the same law, $\Phi_\xi(0)=0$ and $0\in\Med(\xi)$; $0$ is the only nuclear median if in
addition $0$ belongs to the support of $\xi$. Without the last condition there may be others (for $\xi=\pm M_0$ of full
rank, a set of dimension $K(K+1)/2$ when $K=d$).
\item Every nuclear median $C$ of a centred $\xi$ satisfies $\nnuc C\le\E\nnuc{\xi-C}\le\E\nnuc\xi$ (constant $1$ sharp;
without centring the sharp constant is $2$). Hence, in Theorem~\ref{thm:smd}, the limit $W^\dagger$ satisfies
$\nnuc{G^\dagger}^2/(2\Lop)\le F(W^\dagger)-F_\ast\le\nnuc{G^\dagger}^2/(2\muop)\le(\E\nnuc\xi)^2/(2\muop)$.
\item (State-dependent noise.) If the law of the stochastic gradient $\tilde g(W)$ depends on $W$, every zero of the mean
field $W\mapsto\E\,\polar(\tilde g(W))$ is a point where $0$ is a nuclear median of $\tilde g(W)$:
$0\in\E\,\partial\psi(\tilde g(W))$, instead of $\E\,\tilde g(W)=0$.
\item (Finite-sum variance reduction.) For a SAGA or SVRG estimator, at $W_\ast$ with an exact table (or reference
point $W_\ast$) the estimator equals $0$ for every index: the law of the stochastic gradient is a Dirac mass at $0$.
$(W_\ast,\text{exact table})$ is a fixed point of every polar-type step, and the only one: such an estimator has a
fixed point without offset. (No convergence is claimed here; with a fixed length the fixed point is not stable.)
\end{enumerate}
\end{theorem}
\begin{proof}
(a) $\psi_\xi(G)=\E[\nnuc{\xi-(-G)}-\nnuc\xi]$; if $\Phi_\xi(G)=0$ then $0\in\partial\psi_\xi(G)$ by
Theorem~\ref{thm:smooth}, and under the rank condition $\partial\psi_\xi(G)=\{\Phi_\xi(G)\}$. In the first example
$\psi_\xi\ge\frac12\nnuc{\diag(1,-1)}=1=\psi_\xi(\frac12I)+\E\nnuc\xi$ by the triangle inequality, and $\frac12I+\xi$ is
$e_1e_1^\top$ or $e_2e_2^\top$. (b) $\polar$ is odd; $\psi_\xi$ is even and convex, hence minimal at $0$; moreover
$\psi_\xi(G)-\psi_\xi(0)=\frac12\E[\nnuc{\xi+G}+\nnuc{\xi-G}-2\nnuc\xi]$ with an integrand that is nonnegative, and at
least $2\nnuc G-2\nnuc\xi>0$ on the event $\{\nnuc\xi<\nnuc G\}$. (c) $C=\E[C-\xi]$, so
$\nnuc C\le\E\nnuc{\xi-C}\le\E\nnuc\xi$ (Jensen, then minimality of $C$ against $0$); two-point laws show sharpness; the
last line is Proposition~\ref{prop:dual}(c). (d) is Theorem~\ref{thm:smooth} applied to the law of $\tilde g(W)$.
(e) $\omega_j=g_j'(W_\ast)-z_j+\bar z+\lambda W_\ast=\nabla F(W_\ast)=0$ when $z_i=g_i'(W_\ast)$, and $\polar(0)=0$; a
stationary table forces $z_i=g_i'(W)$, hence $\omega_j=\nabla F(W)$ for all $j$.
\end{proof}

\begin{proposition}[batch one on least squares: a least-distances estimator]\label{prop:lad} \textbf{[P]; $K=1$ is classical}
Let $\tilde g(W)=g_j(W)=(Wx_j-y_j)x_j^\top$, $j$ uniform in $\{1,\dots,n\}$ (no ridge inside the polar factor), and
\[
\Lambda(W)=\frac1n\sum_{i=1}^n\frac{\norm{Wx_i-y_i}}{\norm{x_i}} .
\]
Then $\polar(g_j(W))$ is a subgradient at $W$ of $\norm{Wx_j-y_j}/\norm{x_j}$, of Frobenius norm at most $1$, so that
$\E_j\polar(g_j(W))\in\partial\Lambda(W)$, with equality to $\nabla\Lambda(W)$ iff no residual vanishes at $W$. Hence the
batch-one polar step \emph{is} the stochastic subgradient method on $\Lambda$ in the Frobenius geometry: with
$\sum\eta_t=\infty$, $\sum\eta_t^2<\infty$ it converges almost surely to a minimizer of the weighted least-distances
objective $\Lambda$ (least absolute deviations when $K=1$), which is not the least-squares solution in general. Every
zero of the mean field minimizes $\Lambda$; a minimizer at which a residual vanishes need not be a zero of the mean
field (for $K=1$ and generic data with $n>d\ge2$ the mean field has no zero).
\end{proposition}
\begin{proof}
For $r=Wx-y\ne0$, $\polar(rx^\top)=rx^\top/(\norm r\norm x)=\nabla_W\norm{Wx-y}/\norm x$ (Li, Tang and Zou,
arXiv:2602.11557, Lemma~H.1); for $r=0$, $\polar(0)=0\in\partial$. The rest is the standard analysis of the stochastic
subgradient method for a convex Lipschitz function. For $K=1$, $\Lambda$ is polyhedral, its minimizer interpolates $d$
samples for generic data, and the field takes finitely many values $\frac1n\sum_i\varepsilon_ix_i/\norm{x_i}$.
\end{proof}

\begin{proposition}[size of the offset]\label{prop:offset} \textbf{[K] for (a); [F]+[N] for (b)--(c)}
\begin{enumerate}[label=(\alph*)]
\item (Scalar case.) For $K=d=1$ the polar step is sign descent and the limit for batch size $b$ is where the median of
the minibatch gradient vanishes. If the noise is the mean of $b$ i.i.d.\ centred variables with variance $\sigma^2$,
third cumulant $\kappa_3$, finite third absolute moment and a \emph{non-lattice} law, its median is
$-\kappa_3/(6\sigma^2b)+o(1/b)$ (Cornish--Fisher). For a lattice law the median is still $O(1/b)$ but oscillates (for
$X\in\{-1,2\}$ with $\Prob(X=2)=\frac13$, $b$ times the median is $-1,1,0,-1,1,0,\dots$ instead of $-\frac16$). Example:
$f_i(w)=\frac12(w-a_i)^2$, $a_i+1$ exponential: the limit is $\mathrm{median}(\mathrm{Gamma}(b))/b-1\approx-1/(3b)$ and
its objective gap is $\approx1/(18b^2)$, for every schedule with $\sum_t\eta_t=\infty$, $\sum_t\eta_t^2<\infty$.
\item (Matrix case, first order.) Let $\xi_b$ be the mean of $b$ i.i.d.\ centred matrices with covariance operator $C$
and third cumulant tensor $\kappa_3$, and $Z\sim N(0,C)$. Since $\polar$ is homogeneous of degree $0$,
$\Phi_{\xi_b}(0)=\E\polar(\sqrt b\,\xi_b)$, and a one-term Edgeworth expansion gives formally
\[
\Phi_{\xi_b}(0)=\frac{\mathcal T}{\sqrt b}+O(b^{-3/2}),\quad D\Phi_{\xi_b}(0)=\sqrt b\,\mathcal J+\dots,\quad
G_b^\dagger=-\frac1b\,\mathcal J^{-1}\mathcal T+o\Big(\frac1b\Big),
\]
with $\mathcal T=\frac16\E[\polar(Z)\inner{\kappa_3}{\mathrm{He}_3(Z)}]$ and $\mathcal J=D\Phi_Z(0)$ (Proposition~\ref{prop:lowsnr}).
For i.i.d.\ entries of variance $\sigma^2$ and third cumulant $\kappa_3$,
\begin{equation}\label{eq:offset}
G_b^\dagger=\frac{\kappa_3}{6\sigma^2b}\cdot\frac{a_{K,d}}{c_{K,d}}\ \ones\ones^\top+o\Big(\frac1b\Big),\qquad
a_{K,d}=-\E\big[\polar(Z)_{11}(Z_{11}^3-3Z_{11})\big],\quad c_{K,d}=\frac{\E\nnuc Z}{Kd},
\end{equation}
$Z$ standard Gaussian: the scalar offset $\kappa_3/(6\sigma^2b)$ of (a) ($a_{1,1}=c_{1,1}=\sqrt{2/\pi}$) times the factor
$a_{K,d}/c_{K,d}$, in every entry. The objective gap of the limit is of order $1/b^2$. By Gaussian integration by parts
$a_{K,d}=-\E[\nnuc Z\,\mathrm{He}_4(Z_{11})]$, and for $K=1$ one finds $a_{1,d}=3\,\E\chi_d/(d(d+2))$ and
$c_{1,d}=\E\chi_d/d$, so $a_{1,d}/c_{1,d}=3/(d+2)$ exactly. Numerically $a_{K,d}/c_{K,d}$ is of order $1/d$
(Table~\ref{tab:ac}: $d\,a_{K,d}/c_{K,d}$ between $1$ and $2.8$, about $2$ for square matrices): the offset per entry
is smaller than for sign descent by a factor of order $d$, and the nuclear norm of the offset stays between $1$ and
$2.3$ times the scalar offset for the sizes of the table (it decreases for very wide matrices: $3\sqrt d/(d+2)$ for
$K=1$), against $Kd$ times for the $\ell_1$ norm of the offset of sign descent.
\item (Momentum.) For a plain exponential moving average $m_t=\beta m_{t-1}+(1-\beta)\xi_t$ of i.i.d.\ scalar noise and
a frozen iterate, the cumulants of $m$ are $\kappa_n(m)=\kappa_n(1-\beta)^n/(1-\beta^n)$, and the first-order formula of
(a) gives the median $-\frac{\kappa_3}{6\sigma^2}\cdot\frac{1-\beta^2}{1+\beta+\beta^2}$: averaging reduces the offset by
the factor $\frac{1-\beta^2}{1+\beta+\beta^2}\approx\frac23(1-\beta)$, which is larger, by up to $4/3$, than the
``effective batch'' factor $\frac{1-\beta}{1+\beta}$ that governs the variance. With the Nesterov look-ahead used by Muon
(weights $(1-\beta)(1+\beta)$ on the current gradient and $(1-\beta)\beta^{k+1}$ on older ones) the factor is
$(1-\beta)\big[(1+\beta)^3+\frac{\beta^6}{1-\beta^3}\big]\big/\big[(1+\beta)^2+\frac{\beta^4}{1-\beta^2}\big]$ ($0.125$ and
$0.052$ at $\beta=0.9$ and $0.95$, against $0.070$ and $0.034$).
\end{enumerate}
\end{proposition}

\begin{table}[t]
\centering\small
\begin{tabular}{@{}lccccccccccc@{}}
\toprule
$(K,d)$ & $(1,1)$ & $(2,2)$ & $(2,3)$ & $(3,5)$ & $(4,4)$ & $(8,8)$ & $(16,16)$ & $(32,32)$ & $(4,16)$ & $(8,32)$ & $(16,64)$\\
\midrule
$a_{K,d}$ & 0.798 & 0.391 & 0.286 & 0.161 & 0.176 & 0.0737 & 0.0285 & 0.0107 & 0.0389 & 0.0146 & 0.0053\\
$c_{K,d}$ & 0.798 & 0.573 & 0.504 & 0.400 & 0.412 & 0.295 & 0.210 & 0.149 & 0.240 & 0.170 & 0.121\\
$a/c$ & 1 & 0.683 & 0.568 & 0.402 & 0.428 & 0.250 & 0.136 & 0.072 & 0.162 & 0.086 & 0.044\\
$d\cdot a/c$ & 1 & 1.37 & 1.70 & 2.01 & 1.71 & 2.00 & 2.17 & 2.29 & 2.59 & 2.75 & 2.80\\
$\sqrt{Kd}\cdot a/c$ & 1 & 1.37 & 1.39 & 1.56 & 1.71 & 2.00 & 2.17 & 2.29 & 1.30 & 1.38 & 1.40\\
\bottomrule
\end{tabular}
\caption{Constants of the first-order offset \eqref{eq:offset} (Monte Carlo, \texttt{check\_offset\_first\_order.py};
standard errors of $a$ between $1.4\cdot10^{-4}$ and $7.5\cdot10^{-4}$, up to $3\%$ of $a$ for the largest sizes).
Exceptions: the column $(1,1)$ is exact ($\sqrt{2/\pi}$), and the column $(2,3)$ gives the exact values obtained by
quadrature by the referee of the final check ($a_{2,3}=0.28645$, $c_{2,3}=0.50430$; the script gives
$a_{2,3}=0.2869\pm0.0006$); the same quadrature gives $a_{2,2}=0.39108$, $c_{2,2}=0.57290$, in agreement with the
column $(2,2)$. For $K=1$ the script checks $a_{1,d}/c_{1,d}=3/(d+2)$ ($0.5997$ and $0.2982$ for $d=3,8$). Last row:
nuclear norm of the offset in units of the scalar offset $\kappa_3/(6\sigma^2b)$.}\label{tab:ac}
\end{table}

\noindent\emph{Numerical support.} (a): the exact medians are $-0.3069$, $-0.1608$, $-0.0820$, $-0.0413$, $-0.0208$,
$-0.0104$ for $b=1,2,4,8,16,32$. (b): for $2\times3$ matrices with entries $\mathrm{Exp}(1)-1$ ($\sigma=1$, $\kappa_3=2$), \eqref{eq:offset}
predicts $\sqrt b\,\nF{\Phi_{\xi_b}(0)}\to0.234$ and $b\,\nnuc{G_b^\dagger}\to0.464$ ($0.2339$ and $0.4638$ with the exact
constants); the zero of the mean field computed
directly (\texttt{check\_minibatch\_offset\_exact.py}, $1.5\cdot10^6$ samples, exact Gamma batch means) gives
$0.230$--$0.237$ and $0.450$--$0.498$ for $b=2,\dots,256$, and with symmetrized noise the mean field at $0$ vanishes
up to the Monte-Carlo error ($10^{-3}$); the referee of Section~5 derived the same formula independently (predictions
$0.2336$ and $0.4633$ with their Monte-Carlo constants) and
measured $0.230$--$0.241$ and $0.449$--$0.483$ for $b$ up to $1024$ with $3\cdot10^6$ samples; the fixed point is
proportional to the all-ones matrix (cosine above $0.999$ for $b\le64$; $0.998$ and $0.996$ at $b=128$ and $256$, where
the Monte-Carlo error of the root is no longer negligible). \emph{What matters is the skewness of the gradient
noise, not of the data}: in a regression with symmetric inputs independent of the residuals, the per-sample gradient
$\varepsilon x^\top$ is symmetric in law whatever the skewness of $\varepsilon$, and the population offset is zero (an
experiment of the first version, built on such a design, only measured sampling noise; the referee's replication with
Gaussian residuals gives the same values; on the problem of Figure~\ref{fig:noise} without its intercept, the referee
of the final check finds $\nF{W_{\mathrm{LAD}}-W_{\mathrm{LS}}}=0.024$ instead of $0.170$). With an intercept the
gradient noise is skewed:
on a finite-sum least-squares problem with an intercept ($n=2000$, $K=2$, $d=3$, residuals with independent $\frac12(\mathrm{Exp}(1)-1)$ entries; Figure~\ref{fig:noise}, \texttt{make\_paper\_figures.py}), $\nF{W_{\mathrm{LAD}}-W_{\mathrm{LS}}}=0.170$; batch-one polar steps with decreasing lengths end (iterates averaged over the last half of the run and over 32 runs) at distance $2.6\cdot10^{-3}$ from $W_{\mathrm{LAD}}$ and $0.169$ from $W_{\mathrm{LS}}$, and with the SAGA estimator and the same lengths at distance $2.1\cdot10^{-4}$ from $W_{\mathrm{LS}}$. For $b=1, 2, 4, 8, 16, 32$ the zero of the mean field (estimated from $2\times2\cdot10^5$ sampled batches for $b\ge2$, with a spread up to $8.4\cdot10^{-4}$) has $\nnuc{\nabla F}$ equal to $0.168$, $0.120$, $0.069$, $0.036$, $0.019$, $0.0096$ ($b$ times this: $0.168$, $0.240$, $0.276$, $0.285$, $0.303$, $0.308$), stochastic runs with decreasing lengths (same averaging, 64 runs) end at $0.166$, $0.118$, $0.068$, $0.035$, $0.018$, $0.0092$, and the SAGA estimator with the nuclear-scaled length ends below $2.4\cdot10^{-15}$ for every $b$. The referee's population version of such a design gives $b$ times the offset equal to $0.263$, $0.286$, $0.302$, $0.306$, $0.315$ for $b=4,\dots,64$. In this design the per-sample gradients have rank one and the entries of the noise are not independent, so \eqref{eq:offset} does not apply as such: what is observed is the $1/b$ regime, a plateau of $b$ times the offset, reached more slowly than for noise with independent entries. (c): exact medians of the stationary average of $\mathrm{Exp}(1)-1$ noise (referee's computation, hypoexponential
series) are $-0.0466,-0.0227,-0.0112$ for $\beta=0.8,0.9,0.95$; the formula gives $-0.0492,-0.0234,-0.0114$ (ratio
tending to $1$) and the effective-batch heuristic $-0.0370,-0.0175,-0.0085$ (ratio tending to $4/3$).

\begin{figure}[t]
\centering
\includegraphics[width=\textwidth]{figT2_noise}
\caption{Stochastic polar steps on a finite-sum least-squares problem ($n=2000$, $K=2$, $d=3$, intercept, skewed
residuals). (a) Batch one, same decreasing lengths: with the minibatch gradient the averaged iterate goes to the
least-distances solution $W_{\mathrm{LAD}}$ and stays at distance $\nF{W_{\mathrm{LAD}}-W_{\mathrm{LS}}}$ from the
least-squares solution; with the SAGA estimator it goes to $W_{\mathrm{LS}}$. (b) Batch size times the nuclear norm of
the gradient at the limit point: zeros of the mean field (line), stochastic runs with decreasing lengths (dots), SAGA
estimator with the nuclear-scaled length (squares; the gradient norms are at rounding level, below $3\cdot10^{-15}$).
(c) Scalar case with momentum: median of an exponential average of skewed noise against the two factors of
Proposition~\ref{prop:offset}(c). Script: \texttt{make\_paper\_figures.py}.}\label{fig:noise}
\end{figure}

\begin{proposition}[low signal-to-noise regime]\label{prop:lowsnr} \textbf{[K]; the two trace identities are [P]}
Let $\xi=\sigma Z$ with $Z$ a $K\times d$ standard Gaussian matrix, any $K,d\ge1$. Then $\Phi_\xi$ is smooth and odd,
$D\Phi_\xi(G)[\Delta]=\sigma^{-2}\E[\polar(G+\xi)\inner\xi\Delta]$, and
\[
\Phi_\xi(G)=\frac{c_{K,d}}\sigma\,G+O\big(\nF G^3/\sigma^3\big),\qquad c_{K,d}=\frac{\E\nnuc Z}{Kd}\ \in\Big[\frac{\E\chi_{Kd}}{Kd},\ \frac1{\sqrt{\max(K,d)}}\Big].
\]
In the noise-dominated regime the mean polar step is the gradient step with step size $\eta\,c_{K,d}/\sigma$:
\emph{inversely proportional to the noise level}. With $s_1\ge\dots\ge s_K$ the singular values of $Z$ ($K\le d$):
\[
\E\nnuc Z=\E\Big[\sum_{i<j}\frac2{s_i+s_j}+(d-K)\sum_i\frac1{s_i}\Big]\ (K<d),\qquad
\E\nnuc Z=\E\Big[\sum_{i<j}\frac2{s_i+s_j}\Big]+\E\chi_K\ (K=d),
\]
where $\E\chi_K=\sqrt2\,\Gamma(\frac{K+1}2)/\Gamma(\frac K2)$ is the density at $0$ of the smallest singular value of a
$K\times K$ standard Gaussian matrix.
\end{proposition}
\begin{proof}
$\Phi_\xi(G)=\int\polar(y)\phi_\sigma(y-G)\,dy$ with $\polar$ bounded: differentiate the Gaussian density,
$\nabla_G\phi_\sigma(y-G)=\sigma^{-2}(y-G)\phi_\sigma(y-G)$. At $G=0$ the derivative commutes with sign changes and
permutations of rows and columns, so its commutant argument gives a multiple $c'\,\mathrm{Id}$ with
$c'=\tr D\Phi_\xi(0)/(Kd)=\E\inner{\polar(\xi)}{\xi}/(Kd\sigma^2)=\E\nnuc Z/(Kd\sigma)$; $\Phi_\xi$ is odd, so the
remainder is cubic. This is the proof of Xie (arXiv:2609.32861, Proposition~3.1 and Appendix~A) and of Chen et al.\
(arXiv:2610.31728, Proposition~3.1), reproduced by the referee of Section~5. The identities: for $K<d$ the derivative can
also be taken on $\polar$ (the singular set has codimension $d-K+1\ge2$ and $\E[1/s_K]<\infty$), and the trace of
$D\polar(M)$ is $\sum_{i<j}2/(s_i+s_j)+(d-K)\sum_i1/s_i$; for $K=d$ the polar factor jumps by $2u_Kv_K^\top$ across
$\{\det=0\}$ and the Gauss--Green formula adds the surface term $f_K(0)\,\Delta/K^2$, with $f_K(0)=\E\chi_K$ by Edelman's
density of the smallest singular value (SIAM J.\ Matrix Anal.\ Appl.\ 1988).
\end{proof}
\noindent\emph{Numerical support} (\texttt{check\_smoothing\_median.py}, \texttt{\_2.py}, and the referee's quadratures):
$c_{3,5}=0.4003$, $c_{2,8}=0.3369$, $c_{2,2}=0.5729$, $c_{3,3}=0.4727$, $c_{4,4}=0.4122$; for $K=d=2,3,4$ the regular part
alone is $0.2596,0.2954,0.2947$, and measured responses agree with $c_{K,d}$ to three digits. The first version of these
notes presented the singular-value expression as the definition of $c_{K,d}$, with the square case as a formal
computation, and listed the constant as not found; the constant $\E\nnuc Z/(Kd)$, its proof for all shapes and the
reading ``Muon at low signal-to-noise ratio is a gradient step of size $\eta c/\sigma$'' are in Xie and in Chen et al.
The proportional limit ($c_{N,N}\sqrt N\to8/(3\pi)$) and the ``square-root preconditioning'' by a structured noise
covariance are in Paquette et al.\ (arXiv:2605.09552, Theorems 4.1 and K.1). For Gaussian noise $\xi=ZC_R^{1/2}$ the
exact response is $D\Phi_\xi(0)[\Delta]=\Delta A$ with $A=\frac1KC_R^{-1}\E[(\xi^\top\xi)^{1/2}]=(KC_R)^{-1/2}(1+O(d/K))$.

\begin{remark}[three floors, and what acts on each]\label{rem:floors}
For a polar step with stochastic gradients there are three distinct obstructions to precision.
(1) \emph{Length floor}: with a constant spectral length $\eta$, $F-F_\ast\ge\muop\eta^2/8$ at one of any two consecutive
iterates with nonzero gradient (Theorem~\ref{thm:steps}(c)); removed by a nuclear-scaled, clipped or decreasing length,
not by variance reduction.
(2) \emph{Median offset}: for lengths with $\sum_t\eta_t=\infty$, $\sum_t\eta_t^2<\infty$ the limit is a point where
zero is a nuclear median of the stochastic gradient (Theorems~\ref{thm:smd} and~\ref{thm:median}); the set of possible
limits does not depend on the lengths, so no schedule of this class removes the offset. When the law of the stochastic
gradient at $W_\ast$ is symmetric, $W_\ast$ is a possible limit, and the only one if $0$ is in the support of the noise
(for $\xi=\pm I_2$ and $F=\frac12\nF W^2$, $\psi_\xi(wI)=2\max(\abs w,1)-2$, so every $wI$ with $\abs w\le1$ minimizes
$\psi_\xi$; in the simulations of the referee of the final check, runs started at $W_\ast=0$ end at points $wI$ spread
over this segment); the offset is formally of order $1/b$ when the gradient noise
has a nonzero third cumulant; a plain exponential average multiplies it, to first order and for a frozen iterate, by
$(1-\beta^2)/(1+\beta+\beta^2)$; a finite-sum estimator whose noise vanishes at the optimum has a fixed point without
offset. In dimension one, for the regularized Newton--Schulz map with momentum, the corresponding non-convergence for
data with a nonzero third central moment is proved by Do, Dereich and Jentzen (arXiv:2608.04607, Theorem~1.3; their
Proposition~6.2 characterizes limits as zeros of the mean field), with a bound $C(M^{-1}+\gamma_n)$ on the mean-square
error (Theorem~1.2); in their regime the mean field at the minimizer is of order $M^{-3/2}$, not $b^{-1/2}$ as for the
exact polar factor.
(3) \emph{Fluctuations}: of the order of the step length, as for any stochastic method with a constant step.
Experimentally (\texttt{findings.md}, Section~5.1), with constant parameters only the finite-sum estimators combined
with a nuclear-scaled or clipped length reach rounding precision.
\end{remark}

\begin{remark}[attribution]\label{rem:noise}
\emph{Medians.} That sign-based methods drive the median, not the mean, of the stochastic gradient to zero is known:
Chen, Chen, Sun, Wu and Hong (NeurIPS 2020, arXiv:1906.01736); success-probability and symmetric-noise conditions in
Bernstein et al.\ (arXiv:1802.04434, arXiv:1810.05291) and Safaryan and Richt\'arik (arXiv:1905.12938, Lemma~3, with the
ratio third moment over variance); non-convergence of batch-one signSGD for every step-size sequence, when all data
points share a sign pattern, in Karimireddy et al.\ (arXiv:1901.09847, Theorem~I). The normalized step on a location
problem is the Robbins--Monro algorithm for the
geometric median (Cardot, C\'enac and Zitt, arXiv:1101.4316); medians in Banach spaces: Minsker (Bernoulli 2015);
medians of gradients: Schaipp et al.\ (arXiv:2402.12828). For steepest descent with minibatches on separable data, Li,
Tang and Zou show that the limit direction depends on the batch size and that variance reduction removes this
dependence (Theorems 4.7 and 4.10); SignSVRG (Chzhen and Schechtman, arXiv:2305.13187) is the sign analogue.
\emph{Not found}: Theorem~\ref{thm:smd} (the combination ``the method's own noise smooths the nuclear norm, and the
iteration is stochastic mirror descent on the smoothed function in gradient space''), the nuclear-median formulation
for matrices with its rank caveat, Proposition~\ref{prop:lad} for $K\ge2$, the first-order offset \eqref{eq:offset} with
its dimension dependence, and the momentum factor (the two cumulant ratios appear in Do--Dereich--Jentzen,
Proposition~6.11, for another functional). \emph{Other.} $\E\sign(g+\sigma z)=\operatorname{erf}(g/(\sigma\sqrt2))$:
Balles and Hennig (ICML 2018). Deterministic smoothings of the nuclear norm: Mustafi, Mukherjee and Sriperumbudur
(arXiv:2605.23871), Jiang et al. Pethick et al.\ (arXiv:2502.07529) note that the lmo of an unbiased stochastic gradient
is not unbiased. Bounds on the bias of the orthogonalized estimator: Jiang and Xiao (arXiv:2609.36600).
\end{remark}

\section{What this says about accelerating Muon or improving its precision}\label{sec:use}

The statements above are about strongly convex problems with exact structure; their use for deep networks is a set of
hypotheses to test, listed with the result that suggests each.
\begin{enumerate}[label=(\roman*)]
\item \emph{Where the speed is} (known; Section~\ref{sec:ls}). Polar steps are exactly curvature-free when
$E_0^\top E_0$ commutes with $\Sigma$, and the extinction time of the flow is within $\sqrt{\kop}$ of $\nop{E_0}$ in
general (within $1.4$ from generic full-rank starts in all instances tried; up to about $\sqrt d$ for nearly rank-one
errors). The tail is slow: on the invariant ray of the flattest direction the
rate is that of gradient descent with step $1/\lambda_1$ divided by $\tr\Sigma/(2\lambda_1)$, slower as soon as the
stable rank of the data exceeds $2$. A map with bounded gain near zero does not re-excite the dead directions; the
five-step Newton--Schulz map of Muon has such a threshold, about $2\cdot10^{-3}$ of the Frobenius norm (Shulgin et al.,
Appendix~G; Paquette et al., Appendix~C).
\item \emph{Which learning-rate rules are principled.} The nuclear-scaled step is the Polyak step of the dual problem,
and a geometric schedule with ratio $\sqrt{1-1/\kop}$ certifies the same rate when $\muop$, $\Lop$ and a bound on the
initial gap are known (Theorem~\ref{thm:steps}); with a constant length the floor is quadratic in the length.
\item \emph{What limits the final accuracy.} Besides the length floor, the nuclear median: stochastic polar steps
converge to a point where the smoothed nuclear norm of the gradient is minimal, which is $W_\ast$ only if $0$ is a
nuclear median of the noise (Theorem~\ref{thm:smd}). The offset is formally of order $1/b$ when the gradient noise is
skewed, and learning-rate decay does not change it; a plain exponential average reduces it by about
$\frac23(1-\beta)$ to first order; a finite-sum variance-reduced estimator has a fixed point without offset. Without
momentum and with the ridge outside the polar factor, batch-one steps with decreasing lengths converge to the minimizer
of a weighted least-distances objective. For matrices the offset per entry is of order $1/(bd)$, against $1/b$ for sign
descent: orthogonalization averages the skewness.
\item \emph{What the noise does to the step} (known). In the noise-dominated regime the mean polar step is a gradient
step of size $\eta c_{K,d}/\sigma\le\eta/(\sigma\sqrt{\max(K,d)})$: the learning rate that matters is $\eta$ divided by
the noise level, so a batch-size increase by $4$ acts on the drift as a learning-rate increase by $2$.
\item \emph{What cannot be expected.} No certified gain from the spectral geometry, or from any other norm, on
ridge-regularized softmax (Proposition~\ref{prop:struct}(c)); $\kop\ge\min(K,d)$ always, with equality only for
$\frac\mu2\nF\cdot^2$; for least squares with $K\ge d$ the certified full-batch rate is that of the mean per-sample
constant.
\end{enumerate}

\section{Status of the statements and open questions}\label{sec:status}

Table~\ref{tab:status} gives, for every statement, its status and the closest published statements found by the
literature checks.

\begin{table}[!t]
\centering\footnotesize
\begin{tabular}{@{}p{3.4cm}p{1.45cm}p{10.65cm}@{}}
\toprule
Statement & Status & Closest published statements (checks of 12 and 16 Nov.\ 2026); what is added\\
\midrule
Prop.~\ref{prop:var} & P & quantity $\Lop$: Balles et al., Shen et al., Jawanpuria et al., Islamov et al.; duality: Kakade et al.; added: reduction to extreme points for a general operator\\
Prop.~\ref{prop:frame} & P / K & barrier: Sidford--Tian (Lemma A.1), John; added: frame bound; rigidity\textsuperscript{$\dagger$}\\
Cor.~\ref{cor:compare}, Ex.~\ref{ex:sharp} & P & $\ell_\infty$/$\ell_1$ analogues of one side: Balles et al., Nutini et al.; added: two-sided comparison; nuclear-rank form and sharp example\textsuperscript{$\dagger$}\\
Prop.~\ref{prop:struct}(a),(b) & P / K & Shen et al.\ (Kronecker upper bound); Davis--Drusvyatskiy, Jawanpuria et al.\ ($\tr\Sigma$, $K<d$); added: attainment, $\muop$, no averaging gain for $K\ge d$\\
Prop.~\ref{prop:struct}(c),(d) & P & Euclidean metrics, $K=1$: note B1; added: the spectral and preconditioned spectral cases; every norm\textsuperscript{$\dagger$}\\
Prop.~\ref{prop:dual}, \ref{prop:repr} & K & Maddison et al.\ (smooth $k$; Lemma 3.2), Kim et al.; Shulgin et al.\ Eq.~(17), Jawanpuria et al.\ Eq.~(8); added: nonsmooth instance\\
Thm~\ref{thm:steps} & P / K & rate: Balles et al., Gluon; Polyak step: D'Orazio et al.; schedule: Goffin, Yang--He--Li; floors: Jawanpuria et al., Gonon et al., Vaswani et al.; added: identification (a), (b) for general $F$, $\limsup$ bound\\
Prop.~\ref{prop:flow} & K & Leplat et al., Cort\'es, Ma--Wu--E, Bungert et al.\ / Br\'ezis; Kang et al.\ (finite-time convergence of the exact spectral flow in a factorization model); added: spectral instance, \eqref{eq:sandwich}\\
Sec.~\ref{sec:ls}, published part & K & Shulgin et al.\ (Thms 1, 3, Props 9, 10, 12, 17, 18, Cors 3--6, Remarks 1, 4, 7, App.~E.2, G), Vasudeva et al., Wang et al., Gonon et al., Ma et al., Jawanpuria et al.\\
Prop.~\ref{prop:phase1} & P & ingredients: Shulgin et al.\ Props 9, 10, 12, 17, 18; added: frozen factor for every schedule, momentum and the flow, $\tau_{\mathrm{ls}}\ge\tau_\ast$; bounds on $\tau_\ast$ and second-step bound\textsuperscript{$\dagger$}\\
Prop.~\ref{prop:Q} & P & ``if'' part of $Q=I$: Shulgin et al.\ Cor.~4; $\ell_\infty$ analogue of (e): Gruntkowska--Richt\'arik; analogous singular-value formula: Beneventano et al.\ (Thm 1(i)); added: the factor $Q$, the conservation law for $d=2$, the integer counterexample; the angle for $d=2$ and the necessary condition (f)\textsuperscript{$\dagger$}\\
Thm~\ref{thm:mad} & P / K & map, two-cluster rays, tightness at $t=0$: Shulgin et al.\ Cors 4--5, Thm 3, Wang et al.; added: bound for $t\ge1$, converse; instability of the rays\textsuperscript{$\dagger$}\\
Prop.~\ref{prop:ext} & P / K & Vasudeva et al.\ (commuting flow), Gupta et al.\ (path lengths), Bungert et al., Du--Su (single-spike discrete recursion); added: the bound for $d=2$\textsuperscript{$\dagger$}\\
Thm~\ref{thm:smooth} & K & Duchi et al., Abernethy et al., Dubois-Taine et al., Jiang et al.; Xie (noise smooths the polar map)\\
Thm~\ref{thm:smd} & P & not found (combination of the two known facts)\\
Thm~\ref{thm:median}, Prop.~\ref{prop:lad} & P & signs: Chen et al.\ 2020, Safaryan--Richt\'arik, Karimireddy et al.; geometric median: Cardot et al.; Do--Dereich--Jentzen (dimension one); added: matrix formulation, rank caveat, least-distances limit\\
Prop.~\ref{prop:offset} & K / F+N & scalar: Cornish--Fisher, Safaryan--Richt\'arik; cumulant ratios of the momentum: Do--Dereich--Jentzen (Prop.~6.11); added: \eqref{eq:offset}, dimension dependence, momentum factors\\
Prop.~\ref{prop:lowsnr} & K & Xie, Chen et al., Paquette et al.; added: the two trace identities\\
\bottomrule
\end{tabular}
\caption{Status: P proved here, K known, F formal, N numerical. ``Added'' means not found by the literature checks, except the items marked $\dagger$, which were added during the revision (Appendix~\ref{app:review}) and were not searched separately.}\label{tab:status}
\end{table}

\paragraph{Open questions.}
(1) \emph{Extinction time.} For (LS), $\nop{E_0}\le T\le\sqrt{\kop}\nop{E_0}$; from generic full-rank starts the observed
values are between $1.0$ and $1.4$ times $\nop{E_0}$. No constant independent of both $d$ and $\Sigma$ exists, already for rank-one errors (path
lengths of gradient flows, Gupta, Balakrishnan and Ramdas, JMLR 2021: of order $\min\{\sqrt d,\sqrt{\log\kappa}\}$).
For $d=2$, $T\le1.59\,\nop{E_0}$ whatever $\Sigma$ (Proposition~\ref{prop:ext}(d)) and the computed supremum is
$\sqrt2$. Is $T\le\sqrt d\,\nop{E_0}$ in general, or $T\le C\sqrt{\log\kappa(\Sigma)}\,\nop{E_0}$?
(2) \emph{Long-run rate of the line-search polar iteration.} Shulgin et al.\ report its slow-down numerically and
Jawanpuria et al.\ write that the rate is not known in closed form. From generic non-commuting starts the exponent per
step is close to that of the flattest invariant ray, $4\lambda_d/\tr\Sigma$ (Remark~\ref{rem:headtail}), although
these rays are linearly unstable for the commuting map (Theorem~\ref{thm:mad}(c)) and commuting runs do not approach
them. Is the flattest ray attracting for the non-commuting iteration, and why?
(3) The non-commuting iteration after the first phase: the root recursion $R\mapsto\abs{R-\gamma H}$ of Shulgin et al.,
solved there for $d=2$; beyond it, the second-step bound of Proposition~\ref{prop:phase1}(c) is the only general
statement we have.
(4) A proof of the offset expansions of Proposition~\ref{prop:offset}(b)--(c) (Edgeworth expansion against a
discontinuous map), and of the scaling of $a_{K,d}/c_{K,d}$ in $1/d$ beyond $K=1$ (is $d\,a_{K,d}/c_{K,d}$ bounded, with
which limits for square and for wide matrices?).
(5) \emph{State-dependent noise.} For a finite sum the law of the noise depends on $W$ and the mean field is not a
gradient in general (it is a subgradient for batch one on least squares); a Lyapunov function for the joint dynamics of
a SAGA table and a nuclear-scaled step remains open (facts F5 and F6 of \texttt{notes\_theory.md} describe the
obstruction).
(6) Whether the head/tail picture and the median offset are visible on a network: both predictions are quantitative
(nuclear rank of the momentum against the stable rank of the activations; dependence of the final loss on the skewness
of the minibatch gradient, on $b$ and on $\beta$).

\appendix
\section{Numerical checks}\label{app:checks}
All scripts are in \texttt{muon/theory\_checks/} and run with \texttt{.venv/Scripts/python} in seconds to minutes,
except the search \texttt{counterexample\_opnorm.py} and the data of the figures (tens of minutes on a loaded machine;
the figures are redrawn from a cache).
\begin{center}\small
\begin{tabular}{@{}p{6.3cm}p{9.3cm}@{}}
\toprule
script & checks\\
\midrule
\texttt{check\_condition\_numbers.py} & Propositions~\ref{prop:var}--\ref{prop:struct} on random operators (ascent values, frame and rank bounds, Kronecker formulas, least-squares identity)\\
\texttt{check\_dual\_mirror.py} & Proposition~\ref{prop:dual}(c), \eqref{eq:repr}, Theorem~\ref{thm:steps}, flow identities on ridge softmax\\
\texttt{check\_quadratic\_dynamics.py} & Propositions~\ref{prop:phase1}, Theorem~\ref{thm:mad}, commuting dynamics, extinction times\\
\texttt{explore\_rotation\_factor.py}, \texttt{explore\_q1\_sign.py}, \texttt{counterexample\_opnorm.py}, \texttt{counterexample\_opnorm\_integer.py} & Proposition~\ref{prop:Q}: searches, the numerical counterexample of the first version (superseded: its top singular value is double), the integer counterexample\\
\texttt{check\_ls\_additions.py} & additions to Section~\ref{sec:ls} that come from the review: Proposition~\ref{prop:phase1}(a)--(c), Proposition~\ref{prop:Q}(c),(f), Theorem~\ref{thm:mad}(c), Proposition~\ref{prop:ext}(d), the selection identity of the convention paragraph, the tail of the scalar map in Remark~\ref{rem:headtail}\\
\texttt{check\_smoothing\_median.py}, \texttt{check\_smoothing\_median\_2.py} & Theorems~\ref{thm:smooth}, \ref{thm:median}, Propositions~\ref{prop:lad}, \ref{prop:offset}(a),(c), \ref{prop:lowsnr}\\
\texttt{check\_minibatch\_offset\_exact.py}, \texttt{check\_offset\_first\_order.py} & Proposition~\ref{prop:offset}(b): zero of the mean field for exact Gamma batch means; constants of Table~\ref{tab:ac} and the prediction \eqref{eq:offset}\\
\texttt{check\_review\_additions.py} & statements added after the review: Example~\ref{ex:sharp}, Proposition~\ref{prop:struct}(c), the conservation law of Proposition~\ref{prop:Q}, \eqref{eq:sandwich}, Theorem~\ref{thm:smd}, the examples of Theorem~\ref{thm:median}, the lattice example\\
\texttt{make\_paper\_figures.py} & Figures~\ref{fig:ls} and~\ref{fig:noise}, with the runs they show\\
\bottomrule
\end{tabular}
\end{center}

\section{Record of the review of 16 November 2026}\label{app:review}
The first version of these notes was submitted to six independent automated referee passes: three re-derived the
mathematics of Sections~\ref{sec:kappa}--\ref{sec:dual}, of Section~\ref{sec:ls} and of
Sections~\ref{sec:noise}--\ref{sec:use} and tried to break every statement numerically; three searched the literature
for each claim and opened the closest sources (about sixty papers; downloaded copies and scripts are kept under
\texttt{\_scratch/muon\_checks/referee/}). Outcome.

\emph{Mathematics}: no fatal error; corrected here: the inequality
of the former Proposition~2.4(d) (false for $\min(K,d)=1$; replaced by the norm-independent statement), the hidden
hypotheses $K\ge d$ and $\lambda=0$ of the ``no averaging gain'' remark, the case $G_t=0$ and the limsup argument of
Theorem~\ref{thm:steps}(c), the evidence for the sharpness of $\kF/r$, the characterization of stationary points as
nuclear medians (true as an inclusion; equality needs a rank condition), the least-distances statement (same caveat),
the constant of Theorem~\ref{thm:median}(c) ($1$ instead of $2$), the lattice caveat of Proposition~\ref{prop:offset}(a),
several fourth decimals, and one numerical experiment that could not test its claim (symmetric gradient noise). In
Section~\ref{sec:ls} (the last report received): the numerical example of non-monotonicity of the spectral norm had a
double top singular value, so that the formula used to read it did not apply, and is replaced by the integer example;
``the bound is attained'' in Theorem~\ref{thm:mad}(b) became a supremum; the proof of the flow in the commuting case
confused a general cancellation with a consequence of commutation (now the convention paragraph of
Section~\ref{sec:ls}); the thin polar factor is now a stated convention; Remark~\ref{rem:headtail} claimed an alignment
with the flattest direction that is false for commuting starts in exact arithmetic, and compared tails with one
spectrum only.

\emph{Literature}: withdrawn novelty claims on the commuting dynamics, the soft-thresholding flow, the map
$a\mapsto\abs{a-\E_pa}$ and its two-valued rays, the identity $J=\tr\Sigma$, the smoothing identity, the low-noise
constant $c_{K,d}$, the dimension barrier, the dual-space reading for smooth preconditioners and the flow identities;
about forty references added. The mathematical referee of Section~\ref{sec:ls} also found, by reading the appendices of
Shulgin et al., that the factorization behind Proposition~\ref{prop:phase1}, the discrete commuting dynamics, the
two-valued rays and the tightness at $t=0$ were published there, which three ``not found'' sentences of the first
version denied.

\emph{Added during the revision}: Proposition~\ref{prop:struct}(c), Example~\ref{ex:sharp}
and the rigidity statement, \eqref{eq:sandwich}, the conservation law for $d=2$, Theorem~\ref{thm:smd}, the
first-order offset \eqref{eq:offset}, and, from the referee of Section~\ref{sec:ls}, the bounds on $\tau_\ast$, the
second-step bound, items (c) and (f) of Proposition~\ref{prop:Q}, the instability of the rays and
Proposition~\ref{prop:ext}(d). Each of these was derived by one party (the author or a referee) and re-derived or
checked numerically by the other (\texttt{check\_review\_additions.py}, \texttt{check\_ls\_additions.py}); items (c)
of Proposition~\ref{prop:struct}, \eqref{eq:offset} and Theorem~\ref{thm:smd} were derived independently by both.

\emph{Final pass.} The condensed paper (\texttt{paper\_spectral\_muon.tex}) and the statements added during the
revision were then submitted to three further referees: two re-derived the mathematics and recomputed the numbers
(conditioning and noise; dynamics), one checked every citation against the sources. They found no false theorem,
proposition or lemma, and 45 points that are corrected in this version and in the paper: over-statements in the
abstract and the discussion of the paper (the hypothesis of the no-free-lunch statement and the class of step lengths
of the convergence theorem had been dropped in the summaries; ``no offset'' was written for ``a fixed point without
offset''; ``the offset vanishes for symmetric noise'' for ``$W_\ast$ is a possible limit''), two range ends and one
column of Table~\ref{tab:ac}, the generalization of three runs into a rule in Remark~\ref{rem:headtail} (tails of
commuting runs) and the reading of the run of Figure~\ref{fig:ls}(c), the scope of ``all instances tried'' for
extinction times, a uniqueness argument for the flow that was missing in the paper, several attributions (the rate at
every step for $d=2$, the symmetric-noise condition, the effective-batch factor, the hypotheses of two cited theorems),
and the distinction, now marked in Table~\ref{tab:status}, between statements that the literature checks searched and
statements added afterwards. The identity $a_{K,d}=-\E[\nnuc Z\,\mathrm{He}_4(Z_{11})]$ and the exact value
$a_{1,d}/c_{1,d}=3/(d+2)$ of Proposition~\ref{prop:offset}(b) come from that pass. The reports of all passes are kept
under \texttt{\_scratch/muon\_checks/referee/reports/}.

\end{document}
